% Generated mechanically from frozen input p09/ko/src/spaces-groupoids.tex.
% Do not edit this generated file; edit the bound locale source instead.

% OK, start here.
%

\setcounter{chapter}{77}
\chapter{대수공간에서의 Groupoid}



\phantomsection
\label{spaces-groupoids-section-phantom}


\section{서론}
\label{spaces-groupoids-section-introduction}

\noindent
이 장에서는 대수공간에서의 groupoid에 관한 일반론을 다룬다.
Keel과 Mori의 아름다운 논문 \cite{K-M}을 읽기를 권한다.

\medskip\noindent
여기서 말하는 내용의 상당 부분은 groupoid 스킴에 관한 장에서
다룬 내용을 되풀이한 것이다. Groupoids, Section
\ref{groupoids-section-introduction}을 보라.
몫 스택에 관한 논의는 여기서 새로 다룬다.


\section{규약}
\label{spaces-groupoids-section-conventions}

\noindent
모든 스킴은 큰 fppf 사이트 $\Sch_{fppf}$에 들어 있다고
항상 가정한다. 또한 여기서 다루는 모든 환 $A$는
$\Spec(A)$가 이 큰 사이트의 대상과 동형이라는 성질을 가진다.

\medskip\noindent
$S$를 스킴이라 하고 $X$를 $S$ 위의 대수공간이라 하자.
이 장과 다음 장에서는 $X \times_S X$로 $X$와 자기 자신의
곱(즉 $S$ 위 대수공간의 범주에서의 곱)을 나타내며,
$X \times X$라고 쓰지 않는다.

\medskip\noindent
사영 사상의 첨자를 $0$부터 붙이는 규약도 계속 사용한다.
따라서 $\text{pr}_0 : X \times_S Y \to X$ 및
$\text{pr}_1 : X \times_S Y \to Y$이다.




\section{표기}
\label{spaces-groupoids-section-notation}

\noindent
$S$를 스킴이라 하자. 이것을 기저 스킴으로 삼으며 모든
대수공간은 $S$ 위에 놓인다. $B$를 $S$ 위의 대수공간이라
하자. 이것을 기저 대수공간으로 삼으며, 다른 대수공간과
스킴들도 흔히 $B$ 위에 놓인다. $X$가 $B$ 위의 대수공간이라고
하면, 이는 $X$가 $S$ 위의 대수공간이고 구조 사상
$X \to B$가 주어져 있다는 뜻이다. 또한 문자 $T$는
$B$ 위의 {\it 시험} 스킴을 나타내는 데 가급적 남겨 둔다.
즉 $T$는 구조 사상 $T \to B$가 주어진 스킴이다.
이 상황에서 $X(T)$는 $T$-값점들로 이루어진 $X$의 집합을
$B$ {\it 위에서} 나타낸다. 식으로 쓰면
$$
X(T) = \Mor_B(T, X).
$$
마찬가지로 또 다른 대수공간 $Y$가 $B$ 위에 주어지면
$$
X(Y) = \Mor_B(Y, X).
$$
위와 같이 대수공간 $X$, $Y$가 $B$ 위에 있고 사상
$f : X \to Y$가 $B$ 위에 주어졌다고 하자.
$T$를 $B$ 위의 임의의 스킴이라 하면 유도된 집합의 사상
$$
f : X(T) \longrightarrow Y(T)
$$
을 얻으며, 이는 스킴 $T$에 함자적이고 그 스킴은 $B$ 위에 있다.
$f$는 $(\Sch/S)_{fppf}$ 위의 층 사상이고 층 $B$ 위에
있으므로, 이 규칙이 $f$를 결정하며 그 역도 성립한다는 것이
분명하다. 더 일반적으로 올곱들 사이의
사상에도 같은 표기를 쓴다. 예를 들어 $X$, $Y$, $Z$가
$B$ 위의 대수공간이고
$m : X \times_B Y \to Z \times_B Z$가
$B$ 위 대수공간의 사상이면, $m$은 $T$-값점들 사이의
사상들의 모음
$$
X(T) \times Y(T) \longrightarrow Z(T) \times Z(T).
$$
에 대응한다고 생각한다. 이후에도 같은 방식이다.

\medskip\noindent
끝으로 대수공간의 두 사상 $f, g : X \to Y$가 $B$ 위에
주어졌다고 하자. 유도 사상 $f, g : X(T) \to Y(T)$가 모든
스킴 $T$에 대하여 같고 그 스킴이 $B$ 위에 있으면 $f = g$이다.
따라서 사상 $f, g : X(Z) \to Y(Z)$도 모든 세 번째
대수공간 $Z$가 $B$ 위에 있을 때 같다. 그러므로 예를 들어
군 대수공간 $G$가 $B$ 위에 있을 때 그 공리를 확인하려면
$T$-값점에 대한 도식의 가환성을 확인하면 충분하다. 여기서
$T$는 $B$ 위의 스킴이고, 이는 아래 Definition
\ref{spaces-groupoids-definition-group-space}에서 하는 방식이다.




\section{동치관계}
\label{spaces-groupoids-section-equivalence-relations}

\noindent
표기는 Groupoids, Section
\ref{groupoids-section-equivalence-relations}을 보라.

\begin{definition}
\label{spaces-groupoids-definition-equivalence-relation}
Section \ref{spaces-groupoids-section-notation}에서와 같이 $B \to S$라 하자.
$U$를 $B$ 위의 대수공간이라 하자.
\begin{enumerate}
\item $U$의 $B$ 위 {\it 전관계}란 대수공간들의 임의의 사상
$j : R \to U \times_B U$를 말하며, 이 대수공간들은 $B$ 위에 있다.
이 경우 $t = \text{pr}_0 \circ j$ 및
$s = \text{pr}_1 \circ j$로 두므로 $j = (t, s)$이다.
\item $U$의 $B$ 위 {\it 관계}란 대수공간들의 단사 사상
$j : R \to U \times_B U$를 말하며, 이 대수공간들은 $B$ 위에 있다.
\item {\it 전동치관계}란 전관계 $j : R \to U \times_B U$로서,
$j : R(T) \to U(T) \times U(T)$의 상이 모든 스킴 $T$가
$B$ 위에 있을 때 동치관계인 것을 말한다.
\item 대수공간의 사상 $R \to U \times_B U$가 $B$ 위에 있다고 하자.
이것이 {\it $U$의 $B$ 위 동치관계}라는 것은, 모든 $T$가
$B$ 위에 있을 때 $T$-값점들 가운데 $R$에 속하는 것들이
$T$-값점들 가운데 $U$에 속하는 것들의 집합에
동치관계를 정의한다는 것과 동치이다.
\end{enumerate}
\end{definition}

\noindent
달리 말하면, 동치관계란 $j$가 관계인 전동치관계이다.

\begin{lemma}
\label{spaces-groupoids-lemma-restrict-relation}
Section \ref{spaces-groupoids-section-notation}에서와 같이 $B \to S$라 하자.
$U$를 $B$ 위의 대수공간이라 하자.
$j : R \to U \times_B U$를 전관계라 하자.
$g : U' \to U$를 $B$ 위 대수공간의 사상이라 하자.
끝으로
$$
R' = (U' \times_B U')\times_{U \times_B U} R
\xrightarrow{j'}
U' \times_B U'
$$
로 두자. 그러면 $j'$은 $U'$의 전관계이며 $B$ 위에 있다.
$j$가 관계이면 $j'$도 관계이다.
$j$가 전동치관계이면 $j'$도 전동치관계이다.
$j$가 동치관계이면 $j'$도 동치관계이다.
\end{lemma}

\begin{proof}
생략한다.
\end{proof}





\begin{definition}
\label{spaces-groupoids-definition-restrict-relation}
Section \ref{spaces-groupoids-section-notation}에서와 같이 $B \to S$라 하자.
$U$를 $B$ 위의 대수공간이라 하자.
$j : R \to U \times_B U$를 전관계라 하자.
$g : U' \to U$를 $B$ 위 대수공간의 사상이라 하자.
Lemma \ref{spaces-groupoids-lemma-restrict-relation}의 전관계
$j' : R' \to U' \times_B U'$를 전관계 $j$의 $U'$로의
{\it 제한} 또는 {\it 당김}이라 한다.
이 상황에서 때로는 $R' = R|_{U'}$라고 쓴다.
\end{definition}

\begin{lemma}
\label{spaces-groupoids-lemma-pre-equivalence-equivalence-relation-points}
Section \ref{spaces-groupoids-section-notation}에서와 같이 $B \to S$라 하자.
$j : R \to U \times_B U$를 $B$ 위 대수공간의 전관계라 하자.
다음 규칙으로 정의한 $|U|$ 위의 관계를 생각하자.
$$
x \sim y
\Leftrightarrow
\exists\ r \in |R| :
t(r) = x,
s(r) = y.
$$
$j$가 전동치관계이면 이것은 동치관계이다.
\end{lemma}

\begin{proof}
$x \sim y$이고 $y \sim z$라고 하자.
$r \in |R|$ 가운데 $t(r) = x$, $s(r) = y$인 것을 택하고,
$r' \in |R|$ 가운데 $t(r') = y$, $s(r') = z$인 것을 택하자.
체 $K$를 택하여 $r$과 $r'$을 사상
$r, r' : \Spec(K) \to R$로 나타내고
$s \circ r = t \circ r'$이 되게 할 수 있다.
$x = t \circ r$, $y = s \circ r = t \circ r'$,
$z = s \circ r'$로 두면
$x, y, z : \Spec(K) \to U$이다.
구성상 $(x, y) \in j(R(K))$이고
$(y, z) \in j(R(K))$이다. $j$가 전동치관계이므로
$(x, z) \in j(R(K))$도 성립한다.
이는 명백히 $x \sim z$를 함의한다.

\medskip\noindent
$\sim$이 반사적이고 대칭적이라는 증명은 생략한다.
\end{proof}












\section{군 대수공간}
\label{spaces-groupoids-section-group-spaces}

\noindent
표기는 Groupoids, Section
\ref{groupoids-section-group-schemes}을 보라.

\begin{definition}
\label{spaces-groupoids-definition-group-space}
Section \ref{spaces-groupoids-section-notation}에서와 같이 $B \to S$라 하자.
\begin{enumerate}
\item {\it $B$ 위의 군 대수공간}이란 쌍 $(G, m)$을 말한다.
여기서 $G$는 $B$ 위의 대수공간이고
$m : G \times_B G \to G$는 다음 성질을 가지는 $B$ 위
대수공간의 사상이다. 모든 스킴 $T$가 $B$ 위에 있을 때
쌍 $(G(T), m)$은 군이다.
\item {\it 군 대수공간의 사상
$\psi : (G, m) \to (G', m')$가 $B$ 위에 있다}는 것은,
대수공간의 사상 $\psi : G \to G'$가 $B$ 위에 있고 모든
$T/B$에 대하여 유도 사상
$\psi : G(T) \to G'(T)$가 군 준동형인 것을 말한다.
\end{enumerate}
\end{definition}

\noindent
$(G, m)$을 대수공간 $B$ 위의 군 대수공간이라 하자.
Groupoids, Section \ref{groupoids-section-group-schemes}의
논의에 의해 $B$ 위 대수공간의 사상
(항등원) $e : B \to G$와 (역원) $i : G \to G$를 얻으며,
모든 $T$에 대하여 사중항 $(G(T), m, e, i)$는 군의 공리를
만족한다.

\medskip\noindent
$(G, m)$, $(G', m')$을 $B$ 위의 군 대수공간이라 하고
$f : G \to G'$를 $B$ 위 대수공간의 사상이라 하자.
정의에 따라 $f$가 $B$ 위 군 대수공간의 사상인 것은 다음
도식이 가환인 것과 동치이다.
$$
\xymatrix{
G \times_B G \ar[r]_-{f \times f} \ar[d]_m &
G' \times_B G' \ar[d]^m \\
G \ar[r]^f & G'
}
$$

\begin{lemma}
\label{spaces-groupoids-lemma-base-change-group-space}
Section \ref{spaces-groupoids-section-notation}에서와 같이 $B \to S$라 하자.
$(G, m)$을 $B$ 위의 군 대수공간이라 하자.
$B' \to B$를 대수공간의 사상이라 하자.
당김 $(G_{B'}, m_{B'})$은 $B'$ 위의 군 대수공간이다.
\end{lemma}

\begin{proof}
생략한다.
\end{proof}







\section{군 대수공간의 성질}
\label{spaces-groupoids-section-properties-group-spaces}

\noindent
이 절에서는 임의의 기저 위에서 성립하는 군 대수공간의 몇 가지
간단한 성질을 모은다.

\begin{lemma}
\label{spaces-groupoids-lemma-group-scheme-separated}
$S$를 스킴이라 하고 $B$를 $S$ 위의 대수공간이라 하자.
$G$를 $B$ 위의 군 대수공간이라 하자.
$G \to B$가 분리(각각 준분리, 국소 분리)인 것은 항등원 사상
$e : B \to G$가 폐몰입(각각 준콤팩트, 몰입)인 것과 동치이다.
\end{lemma}

\begin{proof}
Morphisms of Spaces, Lemma
\ref{spaces-morphisms-lemma-section-immersion}에 의해,
$e$가 폐몰입(각각 준콤팩트, 몰입)인 것은 $G \to B$가
분리(각각 준분리, 국소 분리)이면 성립함을 상기하자.
역으로 다음 도식을 생각하자.
$$
\xymatrix{
G \ar[r]_-{\Delta_{G/B}} \ar[d] &
G \times_B G \ar[d]^{(g, g') \mapsto m(i(g), g')} \\
B \ar[r]^e & G
}
$$
이 도식이 데카르트라는 것은 대수기하학의 함자적 관점에 관한
한 연습문제이다. 달리 말하면 $\Delta_{G/B}$는 $e$의
밑변환이다. 따라서 $e$가 폐몰입(각각 준콤팩트, 몰입)이면
$\Delta_{G/B}$도 그러하다. Spaces, Lemma
\ref{spaces-lemma-base-change-immersions}
(각각 Morphisms of Spaces, Lemma
\ref{spaces-morphisms-lemma-base-change-quasi-compact},
각각 Spaces, Lemma \ref{spaces-lemma-base-change-immersions})를 보라.
\end{proof}

\begin{lemma}
\label{spaces-groupoids-lemma-group-scheme-unramified-or-lqf}
$S$를 스킴이라 하고 $B$를 $S$ 위의 대수공간이라 하자.
$G$를 $B$ 위의 군 대수공간이라 하자. $G \to B$가
국소 유한형이라고 가정하자. 그러면 $G \to B$가
비분기(각각 국소 준유한)인 것은 $G \to B$가 $e(b)$에서
비분기(각각 준유한)이고 이것이 모든 $b \in |B|$에 대하여
성립하는 것과 동치이다.
\end{lemma}

\begin{proof}
Morphisms of Spaces, Lemma
\ref{spaces-morphisms-lemma-where-unramified}
(각각 Morphisms of Spaces, Lemma
\ref{spaces-morphisms-lemma-base-change-quasi-finite-locus})에 의해,
$U \subset G$인 최대 열린 부분공간 가운데 $U \to B$가
비분기(각각 국소 준유한)가 되는 것이 있으며, $U$의 형성은
밑변환과 가환한다.
따라서 $B = \Spec(k)$가 체의 스펙트럼인 경우로 환원된다.
$g \in G(K)$를 확대 $K/k$에 값을 가지는 점이라 하자.
$g$가 $U$에 속하는지 확인할 때 $K$로 밑변환할 수 있다.
그러므로 다음을 보이면 충분하다.
$$
G \to \Spec(k)\text{가 다음 점에서 비분기이다: }e
\Leftrightarrow
G \to \Spec(k)\text{가 다음 점에서 비분기이다: }g
$$
여기서 $k$-유리점은 $g$이다(준유한성과 $g$, $e$에 대해서도
마찬가지이다). $g$에 의한 평행이동은 $G$의 자기동형이고 이는
$k$ 위에 있으므로
분명하다.
\end{proof}

\begin{lemma}
\label{spaces-groupoids-lemma-open-over-which-unramified-or-lqf}
$S$를 스킴이라 하고 $B$를 $S$ 위의 대수공간이라 하자.
$G$를 $B$ 위의 군 대수공간이라 하자. $G \to B$가
국소 유한형이라고 가정하자.
\begin{enumerate}
\item $U \subset B$인 최대 열린 부분공간으로서 $G_U \to U$가 비분기인 것이
존재하며, $U$의 형성은 밑변환과 가환한다.
\item $U \subset B$인 최대 열린 부분공간으로서
$G_U \to U$가 국소 준유한인 것이 존재하며, $U$의 형성은
밑변환과 가환한다.
\end{enumerate}
\end{lemma}

\begin{proof}
Morphisms of Spaces, Lemma
\ref{spaces-morphisms-lemma-where-unramified}
(각각 Morphisms of Spaces, Lemma
\ref{spaces-morphisms-lemma-base-change-quasi-finite-locus})에 의해,
$W \subset G$인 최대 열린 부분공간 가운데 $W \to B$가
비분기(각각 국소 준유한)가 되는 것이 있다. 또한 $W$의 형성은 밑변환과
가환한다. 두 경우 모두 Lemma
\ref{spaces-groupoids-lemma-group-scheme-unramified-or-lqf}에 의해
$U = e^{-1}(W)$이다.
\end{proof}











\section{군 대수공간의 예}
\label{spaces-groupoids-section-examples-group-spaces}

\noindent
$G \to S$가 기저 스킴 $S$ 위의 군 스킴이면,
밑변환 $G_B$를 임의의 대수공간 $B$에 $S$ 위에서 취했을 때,
이는 Lemma \ref{spaces-groupoids-lemma-base-change-group-space}에 의해 $B$ 위의
군 대수공간이다.
% P09-ERR-0869
아래 예들에서 이를 자주 사용한다.

\begin{example}[곱셈 군 대수공간]
\label{spaces-groupoids-example-multiplicative-group}
Section \ref{spaces-groupoids-section-notation}에서와 같이 $B \to S$라 하자.
임의의 스킴 $T$가 $B$ 위에 있을 때 구조층의 전역 절단의 단위원군
$\Gamma(T, \mathcal{O}_T^*)$를 대응시키는 함자를 생각하자.
이는 군 대수공간
$$
\mathbf{G}_{m, B} = B \times_S \mathbf{G}_{m, S}
$$
으로 표현되고 이 공간은 $B$ 위에 있다. 여기서 $\mathbf{G}_{m, S}$는 $S$ 위의 곱셈군
스킴이다. Groupoids, Example
\ref{groupoids-example-multiplicative-group}을 보라.
\end{example}

\begin{example}[군 대수공간으로서의 1의 근]
\label{spaces-groupoids-example-roots-of-unity}
Section \ref{spaces-groupoids-section-notation}에서와 같이 $B \to S$라 하자.
$n \in \mathbf{N}$이라 하자. 임의의 스킴 $T$가 $B$ 위에 있을 때
$\Gamma(T, \mathcal{O}_T^*)$의 부분군 가운데 $n$차 1의 근들로
이루어진 것을 대응시키는 함자를 생각하자.
이는 군 대수공간
$$
\mu_{n, B} = B \times_S \mu_{n, S}
$$
으로 표현되고 이 공간은 $B$ 위에 있다. 여기서 $\mu_{n, S}$는
$n$차 1의 근들의 군 스킴이고 $S$ 위에 있다. Groupoids, Example
\ref{groupoids-example-roots-of-unity}를 보라.
\end{example}

\begin{example}[덧셈 군 대수공간]
\label{spaces-groupoids-example-additive-group}
Section \ref{spaces-groupoids-section-notation}에서와 같이 $B \to S$라 하자.
임의의 스킴 $T$가 $B$ 위에 있을 때 구조층의 전역 절단군
$\Gamma(T, \mathcal{O}_T)$를 대응시키는 함자를 생각하자.
이는 군 대수공간
$$
\mathbf{G}_{a, B} = B \times_S \mathbf{G}_{a, S}
$$
으로 표현되고 이 공간은 $B$ 위에 있다. 여기서 $\mathbf{G}_{a, S}$는 $S$ 위의 덧셈군
스킴이다. Groupoids, Example
\ref{groupoids-example-additive-group}을 보라.
\end{example}

\begin{example}[일반선형 군 대수공간]
\label{spaces-groupoids-example-general-linear-group}
Section \ref{spaces-groupoids-section-notation}에서와 같이 $B \to S$라 하자.
$n \geq 1$이라 하자. 임의의 스킴 $T$가 $B$ 위에 있을 때 군
$$
\text{GL}_n(\Gamma(T, \mathcal{O}_T))
$$
을 대응시키는 함자를 생각하자. 이는 구조층의 전역 절단 위의
가역 $n \times n$ 행렬들로 이루어진다.
이는 군 대수공간
$$
\text{GL}_{n, B} = B \times_S \text{GL}_{n, S}
$$
으로 표현되고 이 공간은 $B$ 위에 있다.
% P09-ERR-0870
여기서 $\mathbf{G}_{m, S}$는 $S$ 위의 일반선형군 스킴이다.
Groupoids, Example
\ref{groupoids-example-general-linear-group}을 보라.
\end{example}

\begin{example}
\label{spaces-groupoids-example-determinant}
Section \ref{spaces-groupoids-section-notation}에서와 같이 $B \to S$라 하자.
$n \geq 1$이라 하자. 행렬식은 군 대수공간의 사상
$$
\det : \text{GL}_{n, B} \longrightarrow \mathbf{G}_{m, B}
$$
을 정의하고 이 사상은 $B$ 위에 있다. 이는 $S$ 위의 행렬식
사상, 즉 Groupoids, Example \ref{groupoids-example-determinant}의 사상의
밑변환이다.
\end{example}

\begin{example}[상수 군 대수공간]
\label{spaces-groupoids-example-constant-group}
Section \ref{spaces-groupoids-section-notation}에서와 같이 $B \to S$라 하자.
$G$를 추상군이라 하자. 임의의 스킴 $T$가 $B$ 위에 있을 때
국소 상수 사상 $T \to G$의 군을 대응시키는 함자를 생각하자
(여기서 $T$에는 자리스키 위상을, $G$에는 이산 위상을 준다).
이는 군 대수공간
$$
G_B = B \times_S G_S
$$
으로 표현되고 이 공간은 $B$ 위에 있다. 여기서 $G_S$는 Groupoids, Example
\ref{groupoids-example-constant-group}에서 도입한 상수 군 스킴이다.
\end{example}






\section{군 대수공간의 작용}
\label{spaces-groupoids-section-action-group-space}

\noindent
표기는 Groupoids, Section
\ref{groupoids-section-action-group-scheme}을 보라.

\begin{definition}
\label{spaces-groupoids-definition-action-group-space}
Section \ref{spaces-groupoids-section-notation}에서와 같이 $B \to S$라 하자.
$(G, m)$을 $B$ 위의 군 대수공간이라 하자.
$X$를 $B$ 위의 대수공간이라 하자.
\begin{enumerate}
\item {\it $G$의 대수공간 $X/B$ 위 작용}이란 사상
$a : G \times_B X \to X$가 $B$ 위에 있고, 모든 스킴 $T$가
$B$ 위에 있을 때 사상 $a : G(T) \times X(T) \to X(T)$가
$G(T)$-집합의 구조를 $X(T)$에 정의하는 것을 말한다.
\item $X$, $Y$가 $B$ 위의 대수공간이고 각각 $G$의 작용을
갖추었다고 하자. {\it 동변} 사상, 더 정확히는
{\it $G$-동변} 사상 $\psi : X \to Y$란 $B$ 위 대수공간의
사상으로서, 모든 $T$가 $B$ 위에 있을 때
$\psi : X(T) \to Y(T)$가 $G(T)$-집합의 사상인 것을 말한다.
\end{enumerate}
\end{definition}

\noindent
상황 (1)에서는 도식들
\begin{equation}
\label{spaces-groupoids-equation-action}
\xymatrix{
G \times_B G \times_B X \ar[r]_-{1_G \times a} \ar[d]_{m \times 1_X} &
G \times_B X \ar[d]^a \\
G \times_B X \ar[r]^a & X
}
\quad
\xymatrix{
G \times_B X \ar[r]_-a & X \\
X\ar[u]^{e \times 1_X} \ar[ru]_{1_X}
}
\end{equation}
이 가환이라는 뜻이다. 상황 (2)에서는 단지 도식
$$
\xymatrix{
% P09-ERR-0871
G \times_B X \ar[r]_-{\text{id} \times f} \ar[d]_a &
G \times_B Y \ar[d]^a \\
X \ar[r]^f & Y
}
$$
이 가환이라는 뜻이다.

\begin{definition}
\label{spaces-groupoids-definition-free-action}
Definition \ref{spaces-groupoids-definition-action-group-space}에서와 같이
$B \to S$, $G \to B$, $X \to B$라 하자.
$a : G \times_B X \to X$를 $G$의 $X/B$ 위 작용이라 하자.
모든 스킴 $T$가 $B$ 위에 있을 때
$a : G(T) \times X(T) \to X(T)$가 군 $G(T)$의 집합
$X(T)$ 위 자유 작용이면, 이 작용을 {\it 자유롭다}고 한다.
\end{definition}

\begin{lemma}
\label{spaces-groupoids-lemma-free-action}
Definition \ref{spaces-groupoids-definition-free-action}과 같은 상황이라 하자.
% P09-ERR-0872
작용 $a$가 자유로운 것은
$$
G \times_B X \to X \times_B X, \quad (g, x) \mapsto (a(g, x), x)
$$
가 대수공간의 단사 사상인 것과 동치이다.
\end{lemma}

\begin{proof}
정의에서 즉시 따른다.
\end{proof}










\section{주동차 공간}
\label{spaces-groupoids-section-principal-homogeneous}

\noindent
이 절은 Groupoids, Section
\ref{groupoids-section-principal-homogeneous}에 대응한다.
먼저 그 절을 읽기를 권한다.

\begin{definition}
\label{spaces-groupoids-definition-pseudo-torsor}
$S$를 스킴이라 하자. $B$를 $S$ 위의 대수공간이라 하자.
$(G, m)$을 $B$ 위의 군 대수공간이라 하자.
$X$를 $B$ 위의 대수공간이라 하고,
$a : G \times_B X \to X$를 $G$의 $X$ 위 작용이라 하자.
\begin{enumerate}
\item $X$를 {\it 유사 $G$-토서}라 하거나 $X$가
{\it $G$ 아래 형식적 주동차}라고 하는 것은 유도 사상
$G \times_B X \to X \times_B X$,
$(g, x) \mapsto (a(g, x), x)$가 동형일 때이다.
\item 유사 $G$-토서 $X$에 대하여
% P09-ERR-0873
$G$-동변 동형 $G \to X$가 $B$ 위에 존재하고 $G$가
$G$ 위에 왼쪽 곱셈으로 작용하면, 이를 {\it 자명하다}고 한다.
\end{enumerate}
\end{definition}

\noindent
대수공간의 사상 $B' \to B$가 주어졌다고 하자. 그러면
당김 $X_{B'}$은, 원래의 유사 $G$-토서가 $B$ 위에 있을 때,
유사 $G_{B'}$-토서이고 $B'$ 위에 있음이 분명하다.

\begin{lemma}
\label{spaces-groupoids-lemma-characterize-trivial-pseudo-torsors}
Definition \ref{spaces-groupoids-definition-pseudo-torsor}의 상황에서 다음이 성립한다.
\begin{enumerate}
\item 대수공간 $X$가 유사 $G$-토서인 것은 모든 스킴 $T$가
$B$ 위에 있을 때 집합 $X(T)$가 공집합이거나 군 $G(T)$의
$X(T)$ 위 작용이 단순 추이적인 것과 동치이다.
\item 유사 $G$-토서 $X$가 자명한 것은 사상 $X \to B$가
절단을 갖는 것과 동치이다.
\end{enumerate}
\end{lemma}

\begin{proof}
생략한다.
\end{proof}

\begin{definition}
\label{spaces-groupoids-definition-principal-homogeneous-space}
$S$를 스킴이라 하자.
$B$를 $S$ 위의 대수공간이라 하자.
$(G, m)$을 $B$ 위의 군 대수공간이라 하자.
$X$를 유사 $G$-토서라 하고 이것은 $B$ 위에 있다고 하자.
\begin{enumerate}
\item $X$를 {\it 주동차 공간}, 더 정확히는
{\it 주동차 $G$-공간이며 $B$ 위에 있는 것}이라고 하는 것은,
fpqc 덮개\footnote{Groupoids, Definition
\ref{groupoids-definition-principal-homogeneous-space}에서 기본으로 쓰는
토서는 fpqc 덮개 위에서 자명한 유사 토서이다. 한편
% P09-ERR-0874
$G$는 대수공간으로서 군의 층으로 볼 수 있으므로 이미
$G$-토서라는 개념이 있고, 이 개념은
% P09-ERR-0875
fppf-토서에 대응한다. Lemma \ref{spaces-groupoids-lemma-torsor}를 보라.
따라서 fpqc 국소적으로 자명한 유사 토서에는 ``주동차 공간''이라는
말을 쓰며, 이 상황에서는 토서라는 말을 피하려 한다.}
$\{B_i \to B\}_{i \in I}$가 존재하여 각각의
$X_{B_i} \to B_i$가 절단을 가지면(즉 자명한 유사
$G_{B_i}$-토서이면) 성립한다.
\item $\tau \in \{Zariski, \etale, smooth, syntomic, fppf\}$라 하자.
$X$를 {\it $G$-토서}라 하되 그 위상이 $\tau$라고 하거나,
{\it $\tau$ $G$-토서}, 간단히 {\it $\tau$ 토서}라고 하는 것은
$\tau$ 덮개 $\{B_i \to B\}_{i \in I}$가 존재하여 각각의
$X_{B_i} \to B_i$가 절단을 가질 때이다.
\item $X$가 주동차 $G$-공간이고 $B$ 위에 있다고 하자. 이것이
에탈 위상에 대한 토서이면 {\it 준등자명}이라고 한다.
\item $X$가 주동차 $G$-공간이고 $B$ 위에 있다고 하자. 이것이
자리스키 위상에 대한 토서이면 {\it 국소 자명}이라고 한다.
\end{enumerate}
\end{definition}

\noindent
때로는 ``$X$를 $G$-주동차 공간이라 하고 $B$ 위에 두자''라고 말한다.
이는 $X$가 $B$ 위의 대수공간이고 $G$의 작용을 갖추어
$B$ 위의 주동차 공간이 됨을 뜻한다. 이제 두 표기가 모두
적용될 때 앞에서 도입한 표기와 일치함을 보인다.

\begin{lemma}
\label{spaces-groupoids-lemma-torsor}
$S$를 스킴이라 하자.
$(G, m)$을 $S$ 위의 군 대수공간이라 하자.
$X$를 $S$ 위의 대수공간이라 하고,
$a : G \times_S X \to X$를 $G$의 $X$ 위 작용이라 하자.
그러면 Definition \ref{spaces-groupoids-definition-principal-homogeneous-space}의 의미에서
$X$가 $G$-토서이고 그 위상이 $fppf$인 것은 Cohomology on Sites,
Definition \ref{sites-cohomology-definition-torsor}의 의미에서
$X$가 $G$-토서이고 $(\Sch/S)_{fppf}$ 위에 있는 것과 동치이다.
\end{lemma}

\begin{proof}
생략한다.
\end{proof}

\begin{lemma}
\label{spaces-groupoids-lemma-pseudo-torsor-implications}
$S$를 스킴이라 하자. $B$를 $S$ 위의 대수공간이라 하자.
$G$를 $B$ 위의 군 대수공간이라 하자.
$X$를 유사 $G$-토서라 하고 이것은 $B$ 위에 있다고 하자.
% P09-ERR-0876
$G$와 $X$가 $B$ 위에서 국소 유한형이라고 가정하자.
\begin{enumerate}
\item $G \to B$가 비분기이면 $X \to B$도 비분기이다.
\item $G \to B$가 국소 준유한이면 $X \to B$도 국소 준유한이다.
\end{enumerate}
\end{lemma}

\begin{proof}
(1)의 증명. Morphisms of Spaces, Lemma
\ref{spaces-morphisms-lemma-where-unramified}에 의해 $B$가 체의
스펙트럼인 경우로 환원된다. $X$가 공집합이면 결론이 성립한다.
$X$가 공집합이 아니면 체를 확대하여 $X$가 점을 가진다고
가정할 수 있다. 그러면 $G \cong X$이고 결론이 성립한다.

\medskip\noindent
(2)의 증명도 Morphisms of Spaces, Lemma
\ref{spaces-morphisms-lemma-base-change-quasi-finite-locus}를 사용하면
완전히 같은 방식으로 진행된다.
\end{proof}
















\section{동변 준연접층}
\label{spaces-groupoids-section-equivariant}

\noindent
Groupoids, Section \ref{groupoids-section-equivariant}와 비교하라.

\begin{definition}
\label{spaces-groupoids-definition-equivariant-module}
Section \ref{spaces-groupoids-section-notation}에서와 같이 $B \to S$라 하자.
$(G, m)$을 $B$ 위의 군 대수공간이라 하고,
$a : G \times_B X \to X$를 $G$의 작용이라 하자. 여기서
대수공간 $X$는 $B$ 위에 있다.
{\it $G$-동변 준연접 $\mathcal{O}_X$-가군}, 또는 간단히
{\it 동변 준연접 $\mathcal{O}_X$-가군}이란 쌍
$(\mathcal{F}, \alpha)$를 말한다. 여기서 $\mathcal{F}$는 준연접
$\mathcal{O}_X$-가군이고, $\alpha$는 $\mathcal{O}_{G \times_B X}$-가군 사상
$$
\alpha : a^*\mathcal{F} \longrightarrow \text{pr}_1^*\mathcal{F}
$$
이며, $\text{pr}_1 : G \times_B X \to X$는 사영이고 다음을 만족한다.
\begin{enumerate}
\item 도식
% P09-ERR-0877
$$
\xymatrix{
(1_G \times a)^*\text{pr}_2^*\mathcal{F} \ar[r]_-{\text{pr}_{12}^*\alpha} &
\text{pr}_2^*\mathcal{F} \\
(1_G \times a)^*a^*\mathcal{F} \ar[u]^{(1_G \times a)^*\alpha} \ar@{=}[r] &
(m \times 1_X)^*a^*\mathcal{F} \ar[u]_{(m \times 1_X)^*\alpha}
}
$$
% P09-ERR-0878
은 $\mathcal{O}_{G \times_B G \times_B X}$-가군의 범주에서 가환한다.
\item 당김
$$
(e \times 1_X)^*\alpha : \mathcal{F} \longrightarrow \mathcal{F}
$$
은 항등 사상이다.
\end{enumerate}
설명은 Equation (\ref{spaces-groupoids-equation-action})의 관련 도식들과 비교하라.
\end{definition}

\noindent
첫 번째 도식의 가환성은 $(e \times 1_X)^*\alpha$가 $\mathcal{F}$ 위의
멱등 연산자임을 보장한다. 따라서 조건 (2)는 이것이 동형이라는
조건일 뿐이다.

\begin{lemma}
\label{spaces-groupoids-lemma-pullback-equivariant}
Section \ref{spaces-groupoids-section-notation}에서와 같이 $B \to S$라 하자.
$G$를 $B$ 위의 군 대수공간이라 하자.
$f : X \to Y$를 $G$-동변 사상이라 하자. 여기서 두 대수공간은
$B$ 위에 있고 $G$-작용을 갖춘다.
% P09-ERR-0879
당김 $f^*$를
$(\mathcal{F}, \alpha) \mapsto (f^*\mathcal{F}, (1_G \times f)^*\alpha)$
로 정의하면, 준연접 $G$-동변층 가운데 $Y$ 위에 있는 것들의
범주에서 준연접 $G$-동변층 가운데 $X$ 위에 있는 것들의 범주로
가는 함자가 된다.
\end{lemma}

\begin{proof}
생략한다.
\end{proof}





\section{대수공간에서의 groupoid}
\label{spaces-groupoids-section-groupoids}

\noindent
표기는 Groupoids, Section \ref{groupoids-section-groupoids}을 보라.

\begin{definition}
\label{spaces-groupoids-definition-groupoid}
Section \ref{spaces-groupoids-section-notation}에서와 같이 $B \to S$라 하자.
\begin{enumerate}
\item {\it $B$ 위 대수공간에서의 groupoid}란 오중항
$(U, R, s, t, c)$를 말한다. 여기서 $U$와 $R$은 $B$ 위의
대수공간이고, $s, t : R \to U$ 및
$c : R \times_{s, U, t} R \to R$은 다음 성질을 가지는
$B$ 위 대수공간의 사상이다. 임의의 스킴 $T$가 $B$ 위에 있을 때 오중항
$$
(U(T), R(T), s, t, c)
$$
은 groupoid 범주이다.
\item {\it 대수공간에서의 groupoid의 사상
$f : (U, R, s, t, c) \to (U', R', s', t', c')$가 $B$ 위에 있다}는
것은 대수공간의 사상 $f : U \to U'$ 및 $f : R \to R'$가
$B$ 위에 주어지고 다음 성질을 만족하는 것을 말한다. 임의의 스킴
$T$가 $B$ 위에 있을 때 사상 $f$는 groupoid 범주
$(U(T), R(T), s, t, c)$에서 groupoid 범주
$(U'(T), R'(T), s', t', c')$로 가는 함자를 정의한다.
\end{enumerate}
\end{definition}

\noindent
$(U, R, s, t, c)$를 대수공간에서의 groupoid라 하고 이것은 $B$ 위에
있다고 하자. 대수공간의 유일한 사상 $e : U \to R$ 및
$i : R \to R$가 $B$ 위에 존재하여, 모든 스킴 $T$가 $B$ 위에 있을 때
유도 사상
% P09-ERR-0880
$e : U(T) \to R(T)$는 항등이고,
% P09-ERR-0881
$i : R(T) \to R(T)$는 groupoid 범주의 역원이다. 칠중항
$(U, R, s, t, c, e, i)$은 Groupoids, Section
\ref{groupoids-section-groupoids}의 공리 (1), (2)(a), (2)(b), (3)(a),
(3)(b)에 각각 대응하는 가환 도식들을 만족한다. 반대로 이 성질을
가지는 칠중항이 주어지면 오중항 $(U, R, s, t, c)$는 대수공간에서의
groupoid이고 $B$ 위에 있다. $i$는 동형이고 $e$는 $s$와 $t$ 모두의
절단임에 유의하라. 또한 $B$ 위 대수공간에서의 groupoid가 주어지면
다음과 같이 쓴다.
$$
j = (t, s) : R \longrightarrow U \times_B U
$$
이는 위의 Section \ref{spaces-groupoids-section-equivalence-relations}에서 정한 관례와
양립한다. 항등원과 역원의 존재를 강조하려고 때로는
``$(U, R, s, t, c, e, i)$를 대수공간에서의 groupoid라 하고
$B$ 위에 두자''라고 말한다.

\begin{lemma}
\label{spaces-groupoids-lemma-groupoid-pre-equivalence}
Section \ref{spaces-groupoids-section-notation}에서와 같이 $B \to S$라 하자.
대수공간에서의 groupoid $(U, R, s, t, c)$가 $B$ 위에 주어지면,
사상 $j : R \to U \times_B U$는 전동치관계이다.
\end{lemma}

\begin{proof}
생략한다. 정의에 관한 좋은 연습문제이다.
\end{proof}

\begin{lemma}
\label{spaces-groupoids-lemma-equivalence-groupoid}
Section \ref{spaces-groupoids-section-notation}에서와 같이 $B \to S$라 하자.
동치관계 $j : R \to U \times_B U$가 $B$ 위에 주어지면, 이를
대수공간에서의 groupoid $(U, R, s, t, c)$로 확장하고 $B$ 위에
두는 방법이 유일하게 존재한다.
\end{lemma}

\begin{proof}
생략한다. 정의에 관한 좋은 연습문제이다.
\end{proof}

\begin{lemma}
\label{spaces-groupoids-lemma-diagram}
Section \ref{spaces-groupoids-section-notation}에서와 같이 $B \to S$라 하자.
$(U, R, s, t, c)$를 대수공간에서의 groupoid라 하고 $B$ 위에 두자.
가환 도식
$$
\xymatrix{
& U & \\
R \ar[d]_s \ar[ru]^t &
R \times_{s, U, t} R
\ar[l]^-{\text{pr}_0} \ar[d]^{\text{pr}_1} \ar[r]_-c &
R \ar[d]^s \ar[lu]_t \\
U & R \ar[l]_t \ar[r]^s & U
}
$$
에서 아래쪽 두 정사각형은 올곱 정사각형이다. 또한 위쪽 삼각형은
(실제로는 정사각형인데) 역시 데카르트이다.
\end{lemma}

\begin{proof}
생략한다. 이는 정의와 대수기하학의 함자적 관점에 관한 연습문제이다.
\end{proof}

\begin{lemma}
\label{spaces-groupoids-lemma-diagram-pull}
Section \ref{spaces-groupoids-section-notation}에서와 같이 $B \to S$라 하자.
$(U, R, s, t, c, e, i)$를 대수공간에서의 groupoid라 하고 $B$ 위에 두자.
도식
\begin{equation}
\label{spaces-groupoids-equation-pull}
\xymatrix{
R \times_{t, U, t} R
\ar@<1ex>[r]^-{\text{pr}_1} \ar@<-1ex>[r]_-{\text{pr}_0}
\ar[d]_{\text{pr}_0 \times c \circ (i, 1)} &
R \ar[r]^t \ar[d]^{\text{id}_R} &
U \ar[d]^{\text{id}_U} \\
R \times_{s, U, t} R
\ar@<1ex>[r]^-c \ar@<-1ex>[r]_-{\text{pr}_0} \ar[d]_{\text{pr}_1} &
R \ar[r]^t \ar[d]^s &
U \\
R \ar@<1ex>[r]^s \ar@<-1ex>[r]_t &
U
}
\end{equation}
은 가환한다. 위쪽 두 행은 주어진 수직 사상들을 통해 동형이다.
왼쪽 아래의 두 정사각형은 데카르트이다.
\end{lemma}

\begin{proof}
도식의 가환성은 groupoid의 공리에서 따른다. groupoid의 말로 하면,
왼쪽 위 수직 화살표는 같은 표적을 가지는 사상의 쌍
$(\alpha, \beta)$에 사상의 쌍
$(\alpha, \alpha^{-1} \circ \beta)$를 대응시킨다. 모든 groupoid에서
이는
$\text{Arrows} \times_{t, \text{Ob}, t} \text{Arrows}$와
$\text{Arrows} \times_{s, \text{Ob}, t} \text{Arrows}$ 사이의
전단사를 정의한다. 따라서 보조정리의 두 번째 주장이 따른다.
마지막 주장은 Lemma \ref{spaces-groupoids-lemma-diagram}에서 따른다.
\end{proof}

\begin{lemma}
\label{spaces-groupoids-lemma-base-change-groupoid}
Section \ref{spaces-groupoids-section-notation}에서와 같이 $B \to S$라 하자.
$(U, R, s, t, c)$를 대수공간에서의 groupoid라 하고 $B$ 위에 두자.
$B' \to B$를 대수공간의 사상이라 하자. 그러면 밑변환
$U' = B' \times_B U$, $R' = B' \times_B R$에 $s'$, $t'$, $c'$를
갖추되 이들이 사상 $s, t, c$의 밑변환이면, 대수공간에서의 groupoid
$(U', R', s', t', c')$가 되고 이것은 $B'$ 위에 있다. 또한 사영들은
대수공간에서의 groupoid의 사상
$(U', R', s', t', c') \to (U, R, s, t, c)$를 정의하며 이것은
$B$ 위에 있다.
\end{lemma}

\begin{proof}
생략한다. 힌트:
$R' \times_{s', U', t'} R' = B' \times_B (R \times_{s, U, t} R)$.
\end{proof}





\section{Groupoid 위의 준연접층}
\label{spaces-groupoids-section-groupoids-quasi-coherent}

\noindent
Groupoids, Section
\ref{groupoids-section-groupoids-quasi-coherent}와 비교하라.

\begin{definition}
\label{spaces-groupoids-definition-groupoid-module}
Section \ref{spaces-groupoids-section-notation}에서와 같이 $B \to S$라 하자.
$(U, R, s, t, c)$를 대수공간에서의 groupoid라 하고 $B$ 위에 두자.
{\it $(U, R, s, t, c)$ 위의 준연접 가군}이란 쌍
$(\mathcal{F}, \alpha)$를 말한다. 여기서 $\mathcal{F}$는 준연접
$\mathcal{O}_U$-가군이고, $\alpha$는 $\mathcal{O}_R$-가군 사상
$$
\alpha : t^*\mathcal{F} \longrightarrow s^*\mathcal{F}
$$
이며 다음을 만족한다.
\begin{enumerate}
\item 도식
$$
\xymatrix{
& \text{pr}_1^*t^*\mathcal{F} \ar[r]_-{\text{pr}_1^*\alpha} &
\text{pr}_1^*s^*\mathcal{F} \ar@{=}[rd] & \\
\text{pr}_0^*s^*\mathcal{F} \ar@{=}[ru] & & & c^*s^*\mathcal{F} \\
& \text{pr}_0^*t^*\mathcal{F} \ar[lu]^{\text{pr}_0^*\alpha} \ar@{=}[r] &
c^*t^*\mathcal{F} \ar[ru]_{c^*\alpha}
}
$$
% P09-ERR-0882
은 $\mathcal{O}_{R \times_{s, U, t} R}$-가군의 범주에서 가환한다.
\item 당김
$$
e^*\alpha : \mathcal{F} \longrightarrow \mathcal{F}
$$
은 항등 사상이다.
\end{enumerate}
Lemma \ref{spaces-groupoids-lemma-diagram}의 가환 도식들과 비교하라.
\end{definition}

\noindent
첫 번째 도식의 가환성은 연산자 $e^*\alpha$가 멱등임을 강제한다.
따라서 두 번째 조건은 $e^*\alpha$가 동형이라는 말로 다시 쓸 수 있다.
실제로 이 조건은 $\alpha$가 동형임을 함의한다.

\begin{lemma}
\label{spaces-groupoids-lemma-isomorphism}
Section \ref{spaces-groupoids-section-notation}에서와 같이 $B \to S$라 하자.
$(U, R, s, t, c)$를 대수공간에서의 groupoid라 하고 $B$ 위에 두자.
$(\mathcal{F}, \alpha)$가 $(U, R, s, t, c)$ 위의 준연접 가군이면
$\alpha$는 동형이다.
\end{lemma}

\begin{proof}
Definition \ref{spaces-groupoids-definition-groupoid-module}의 가환 도식을 사상
$(i, 1) : R \to R \times_{s, U, t} R$로 당겨오자. 그러면
$i^*\alpha \circ \alpha = s^*e^*\alpha$임을 알 수 있다. 사상
$(1, i)$로 당겨오면 관계
$\alpha \circ i^*\alpha = t^*e^*\alpha$를 얻는다. 두 번째 가정에
의해 이 사상들은 항등이다. 따라서 $i^*\alpha$는 $\alpha$의 역이다.
\end{proof}

\begin{lemma}
\label{spaces-groupoids-lemma-pullback}
Section \ref{spaces-groupoids-section-notation}에서와 같이 $B \to S$라 하자. 사상
$f : (U, R, s, t, c) \to (U', R', s', t', c')$를 생각하자. 이는
% P09-ERR-0883
대수공간에서의 groupoid의 사상이고 $B$ 위에 있다. 다음으로 정의한
% P09-ERR-0884
당김 $f^*$
$$
(\mathcal{F}, \alpha) \mapsto (f^*\mathcal{F}, f^*\alpha)
$$
은 $(U', R', s', t', c')$ 위의 준연접층의 범주에서
$(U, R, s, t, c)$ 위의 준연접층의 범주로 가는 함자를 정의한다.
\end{lemma}

\begin{proof}
생략한다.
\end{proof}

\begin{lemma}
\label{spaces-groupoids-lemma-pushforward}
Section \ref{spaces-groupoids-section-notation}에서와 같이 $B \to S$라 하자. 사상
$f : (U, R, s, t, c) \to (U', R', s', t', c')$를 생각하자. 이는
대수공간에서의 groupoid의 사상이고 $B$ 위에 있다. 다음을 가정하자.
\begin{enumerate}
\item $f : U \to U'$는 준콤팩트이고 준분리이다.
\item 정사각형
$$
\xymatrix{
R \ar[d]_t \ar[r]_f & R' \ar[d]^{t'} \\
U \ar[r]^f & U'
}
$$
은 데카르트이다.
\item $s'$와 $t'$는 평탄하다.
\end{enumerate}
% P09-ERR-0885
다음으로 정의한 앞보냄 $f_*$
$$
(\mathcal{F}, \alpha) \mapsto (f_*\mathcal{F}, f_*\alpha)
$$
은 $(U, R, s, t, c)$ 위의 준연접층의 범주에서
$(U', R', s', t', c')$ 위의 준연접층의 범주로 가는 함자를 정의하며,
Lemma \ref{spaces-groupoids-lemma-pullback}에서 정의한 당김의 오른쪽 수반이다.
\end{lemma}

\begin{proof}
$U \to U'$가 준콤팩트이고 준분리이므로 $f_*$는 준연접층을
준연접층으로 보낸다(Morphisms of Spaces, Lemma
\ref{spaces-morphisms-lemma-pushforward}). 또한 정사각형들
$$
\vcenter{
\xymatrix{
R \ar[d]_t \ar[r]_f & R' \ar[d]^{t'} \\
U \ar[r]^f & U'
}
}
\quad\text{및}\quad
\vcenter{
\xymatrix{
R \ar[d]_s \ar[r]_f & R' \ar[d]^{s'} \\
U \ar[r]^f & U'
}
}
$$
이 데카르트이므로
$(t')^*f_*\mathcal{F} = f_*t^*\mathcal{F}$ 및
$(s')^*f_*\mathcal{F} = f_*s^*\mathcal{F}$를 얻는다. Cohomology of
Spaces, Lemma
\ref{spaces-cohomology-lemma-flat-base-change-cohomology}를 보라.
따라서 $f_*\alpha$를 사상
$(t')^*f_*\mathcal{F} \to (s')^*f_*\mathcal{F}$로 생각하는 것은
타당하다. 비슷한 논증으로 $f_*\alpha$가 코사이클 조건을 만족함을
알 수 있다. 환 달린 공간 위 가군의 당김과 앞보냄은 서로 수반이므로
이 함자는 당김 함자와 수반이다. 몇 가지 세부사항은 생략한다.
\end{proof}


\begin{lemma}
\label{spaces-groupoids-lemma-colimits}
Section \ref{spaces-groupoids-section-notation}에서와 같이 $B \to S$라 하자.
$(U, R, s, t, c)$를 대수공간에서의 groupoid라 하고 $B$ 위에 두자.
$(U, R, s, t, c)$ 위의 준연접 가군의 범주는 여극한을 가진다.
\end{lemma}

\begin{proof}
$i \mapsto (\mathcal{F}_i, \alpha_i)$를 도식이라 하고 그 지표 범주를
$\mathcal{I}$라 하자. 여극한
$\mathcal{F} = \colim \mathcal{F}_i$를 만들 수 있으며, 이는 $U$ 위의
준연접층이다. Properties of Spaces, Lemma
\ref{spaces-properties-lemma-properties-quasi-coherent}를 보라. 여극한은
당김과 가환하므로 $s^*\mathcal{F} = \colim s^*\mathcal{F}_i$이고,
마찬가지로 $t^*\mathcal{F} = \colim t^*\mathcal{F}_i$이다. 따라서
$\alpha = \colim \alpha_i$로 둘 수 있다. $(\mathcal{F}, \alpha)$가
$(U, R, s, t, c)$ 위의 준연접 가군의 범주에서 이 도식의 여극한이라는
증명은 생략한다.
\end{proof}

\begin{lemma}
\label{spaces-groupoids-lemma-abelian}
Section \ref{spaces-groupoids-section-notation}에서와 같이 $B \to S$라 하자.
$(U, R, s, t, c)$를 대수공간에서의 groupoid라 하고 $B$ 위에 두자.
$s$, $t$가 평탄하면 $(U, R, s, t, c)$ 위의 준연접 가군의 범주는
아벨 범주이다.
\end{lemma}

\begin{proof}
$\varphi : (\mathcal{F}, \alpha) \to (\mathcal{G}, \beta)$를
$(U, R, s, t, c)$ 위의 준연접 가군의 준동형이라 하자. $s$가
평탄하므로
$$
0 \to s^*\Ker(\varphi)
\to s^*\mathcal{F} \to s^*\mathcal{G} \to s^*\Coker(\varphi) \to 0
$$
은 완전하고 $t$를 따른 당김도 마찬가지이다. 따라서 $\alpha$와
$\beta$는 코사이클 조건을 만족하는 동형
$\kappa : t^*\Ker(\varphi) \to s^*\Ker(\varphi)$ 및
$\lambda : t^*\Coker(\varphi) \to s^*\Coker(\varphi)$를 유도한다.
이제 $(\Ker(\varphi), \kappa)$와 $(\Coker(\varphi), \lambda)$가
$(U, R, s, t, c)$ 위의 준연접 가군 범주에서 각각 핵과 여핵임을
확인하는 것은 곧바르다. 또한 $\Coim(\varphi) = \Im(\varphi)$라는
조건은 $U$ 위에서 성립하므로 여기서도 성립한다.
\end{proof}










\section{준연접 가군의 여극한}
\label{spaces-groupoids-section-colimits}

\noindent
이 절은 Groupoids, Section \ref{groupoids-section-colimits}에 대응한다.

\begin{lemma}
\label{spaces-groupoids-lemma-construct-quasi-coherent}
Section \ref{spaces-groupoids-section-notation}에서와 같이 $B \to S$라 하자.
$(U, R, s, t, c)$를 대수공간에서의 groupoid라 하고 $B$ 위에 두자.
$s, t$가 평탄하고 준콤팩트이며 준분리라고 가정하자. 임의의 준연접
가군 $\mathcal{G}$가 $U$ 위에 있을 때 표준 동형
$\alpha : t^*s_*t^*\mathcal{G} \to s^*s_*t^*\mathcal{G}$가 존재하며,
이는 $(s_*t^*\mathcal{G}, \alpha)$를 $(U, R, s, t, c)$ 위의 준연접
가군으로 만든다. 이 구성은 함자
$$
\QCoh(\mathcal{O}_U) \longrightarrow \QCoh(U, R, s, t, c)
$$
를 정의하며, 이 함자는 망각 함자
$(\mathcal{F}, \beta) \mapsto \mathcal{F}$의 오른쪽 수반이다.
\end{lemma}

\begin{proof}
준콤팩트이고 준분리인 사상을 따른 준연접 가군의 앞보냄은
준연접이다. Morphisms of Spaces, Lemma
\ref{spaces-morphisms-lemma-pushforward}를 보라. 따라서
$s_*t^*\mathcal{G}$는 준연접이다. Lemma \ref{spaces-groupoids-lemma-diagram}의 표기를
쓰면
$$
t^*s_*t^*\mathcal{G} =
\text{pr}_{1, *}\text{pr}_0^*t^*\mathcal{G} =
\text{pr}_{1, *}c^*t^*\mathcal{G} =
s^*s_*t^*\mathcal{G}
$$
이다. 가운데 등식은 $t \circ c = t \circ \text{pr}_0$가 사상
$R \times_{s, U, t} R \to U$ 사이에서 성립하기 때문이고, 첫째 및 마지막 등식은
이 단계에서 밑변환과 앞보냄이 가환함을 알고 있기 때문이다.
Cohomology of Spaces, Lemma
\ref{spaces-cohomology-lemma-flat-base-change-cohomology}를 보라.

\medskip\noindent
Definition \ref{spaces-groupoids-definition-groupoid-module}의 $\alpha$에 대한 코사이클
조건과 수반성을 확인하기 위해 구성
$\mathcal{G} \mapsto (s_*t^*\mathcal{G}, \alpha)$를 다른 방식으로
설명하자.
% P09-ERR-0886
groupoid 스킴
% P09-ERR-0888
$(R, R \times_{t, U, t} R, \text{pr}_0, \text{pr}_1, \text{pr}_{02})$를
생각하자. 이것은 동치관계 $R \times_{t, U, t} R$에 결부되며,
그 관계는 $R$ 위에 있다. Lemma \ref{spaces-groupoids-lemma-equivalence-groupoid}를 보라. 사상
$$
f :
(R, R \times_{t, U, t} R, \text{pr}_1, \text{pr}_0, \text{pr}_{02})
\longrightarrow
(U, R, s, t, c)
$$
% P09-ERR-0887
이 존재한다. 이는 groupoid 스킴의 사상으로서 $s : R \to U$와
$R \times_{t, U, t} R \to R$로 주어지며, 뒤의 사상은
$(r_0, r_1) \mapsto r_0^{-1} \circ r_1$이다. 필요한 도식들의 가환성
확인은 생략한다. $t, s : R \to U$는 준콤팩트이고 준분리이며
평탄하고, 데카르트 정사각형
$$
\xymatrix{
R \times_{t, U, t} R \ar[d]_{\text{pr}_0}
\ar[rr]_-{(r_0, r_1) \mapsto r_0^{-1} \circ r_1} & & R \ar[d]^t \\
R \ar[rr]^s & & U
}
$$
이 있기 때문에, Lemma \ref{spaces-groupoids-lemma-diagram-pull}에 의해
Lemma \ref{spaces-groupoids-lemma-pushforward}를 $f$에 적용할 수 있다. 따라서 $f$를
따른 준연접 가군의 앞보냄과 당김은 수반 함자이다. 증명을 끝내기 위해
이 함자들을 위에서 설명한 함자들과 동일시하겠다. 이를 위해
$$
t^* :
\QCoh(\mathcal{O}_U)
\longrightarrow
\QCoh(R, R \times_{t, U, t} R, \text{pr}_1, \text{pr}_0, \text{pr}_{02})
$$
는 준연접층의 하강 이론에 의해 동치임에 유의하자. 실제로
$\{t : R \to U\}$는 fpqc 덮개이다. Descent on Spaces, Proposition
\ref{spaces-descent-proposition-fpqc-descent-quasi-coherent}를 보라.

\medskip\noindent
$f$를 따른 앞보냄 앞에 동치 $t^*$를 합성하면 $\mathcal{G}$를
$(s_*t^*\mathcal{G}, \alpha)$로 보낸다. 이 방식으로 얻은 동형
$\alpha$가 위에서 구성한 것과 같다는 확인은 생략한다.

\medskip\noindent
$f$를 따른 당김 뒤에 동치 $t^*$의 역을 합성하면
$(\mathcal{F}, \beta)$를 $\{t : R \to U\}$에 관한 가군
$s^*\mathcal{F}$의 하강으로 보낸다. 이 가군에는 하강 자료 $\gamma$를
갖추고, 이 자료는 $R \times_{t, U, t} R$ 위에 있으며 이는
$\beta$를 $(r_0, r_1) \mapsto r_0^{-1} \circ r_1$로 당겨온 것이다.
동형 $\beta : t^*\mathcal{F} \to s^*\mathcal{F}$를 생각하자.
$t^*\mathcal{F}$ 위의 표준 하강 자료로서 $\{t : R \to U\}$에 관한 것
(Descent on Spaces, Definition
\ref{spaces-descent-definition-descent-datum-effective-quasi-coherent})는
$\beta$를 통해 사상
$$
\text{pr}_0^*s^*\mathcal{F}
\xrightarrow{\text{pr}_0^*\beta^{-1}}
\text{pr}_0^*t^*\mathcal{F}
\xrightarrow{can}
\text{pr}_1^*t^*\mathcal{F}
\xrightarrow{\text{pr}_1^*\beta}
\text{pr}_1^*s^*\mathcal{F}
$$
으로 옮겨진다. $\beta$가 코사이클 조건을 만족하므로, 이는
$\beta$를 $(r_0, r_1) \mapsto r_0^{-1} \circ r_1$로 당겨온 것과
같다. 이를 보려면 Definition \ref{spaces-groupoids-definition-groupoid-module}의 실제
코사이클 관계를 취해 사상
$(\text{pr}_0, c \circ (i, 1)) : R \times_{t, U, t} R \to R \times_{s, U, t} R$
로 당겨오라. 이 사상은 Lemma \ref{spaces-groupoids-lemma-diagram-pull}의 가환 도식에도
등장한다. 따라서 $(s^*\mathcal{F}, \gamma)$는
$(t^*\mathcal{F}, can)$과 동형이다. 모두 합치면, $f$를 따른 당김 뒤에
동치 $t^*$의 역을 합성한 것은 망각 함자
$(\mathcal{F}, \beta) \mapsto \mathcal{F}$와 동형이다.
\end{proof}

\begin{remark}
\label{spaces-groupoids-remark-adjunction-map}
Lemma \ref{spaces-groupoids-lemma-construct-quasi-coherent}의 상황에서
$$
F : \QCoh(U, R, s, t, c) \to \QCoh(\mathcal{O}_U),\quad
(\mathcal{F}, \beta) \mapsto \mathcal{F}
$$
를 망각 함자라 하고,
$$
G : \QCoh(\mathcal{O}_U) \to \QCoh(U, R, s, t, c),\quad
\mathcal{G} \mapsto (s_*t^*\mathcal{G}, \alpha)
$$
를 보조정리에서 구성한 오른쪽 수반이라 하자. 수반의 단위원
$\eta : \text{id} \to G \circ F$를 $(\mathcal{F}, \beta)$에서
계산하면 사상
$$
\mathcal{F} \to s_*s^*\mathcal{F} \xrightarrow{\beta^{-1}} s_*t^*\mathcal{F}
$$
으로 주어진다. 확인은 생략한다.
\end{remark}

\begin{lemma}
\label{spaces-groupoids-lemma-push-pull}
$S$를 스킴이라 하자.
$f : Y \to X$를 $S$ 위 대수공간의 사상이라 하자.
$\mathcal{F}$를 준연접 $\mathcal{O}_X$-가군이라 하고,
$\mathcal{G}$를 준연접 $\mathcal{O}_Y$-가군이라 하며,
$\varphi : \mathcal{G} \to f^*\mathcal{F}$를 가군 사상이라 하자.
다음을 가정하자.
\begin{enumerate}
\item $\varphi$는 단사이다.
\item $f$는 준콤팩트이고 준분리이며 평탄하고 전사이다.
\item $X$, $Y$는 국소 뇌터이고,
\item $\mathcal{G}$는 연접 $\mathcal{O}_Y$-가군이다.
\end{enumerate}
올곱으로 정의한 $\mathcal{F} \cap f_*\mathcal{G}$
$$
\xymatrix{
\mathcal{F} \ar[r] & f_*f^*\mathcal{F} \\
\mathcal{F} \cap f_*\mathcal{G} \ar[u] \ar[r] &
f_*\mathcal{G} \ar[u]
}
$$
은 연접 $\mathcal{O}_X$-가군이다.
\end{lemma}

\begin{proof}
Cohomology of Spaces, Lemma
\ref{spaces-cohomology-lemma-coherent-Noetherian}의 연접 가군의 특성화와,
연접 가군들이 $\QCoh(\mathcal{O}_X)$의 Serre 부분범주를 이룬다는
사실을 자유롭게 쓰겠다. 후자는 Cohomology of Spaces, Lemma
\ref{spaces-cohomology-lemma-coherent-Noetherian-quasi-coherent-sub-quotient}을
보라. $f$가 절단 $\sigma$를 가지면
$\mathcal{F} \cap f_*\mathcal{G}$는
$\sigma^*\mathcal{G} \to \sigma^*f^*\mathcal{F} = \mathcal{F}$의 상에
포함되므로 연접이다. 일반적으로 $\mathcal{F} \cap f_*\mathcal{G}$가
연접임을 보이기 위해서는
% P09-ERR-0889
$f^*(\mathcal{F} \cap f_*\mathcal{G})$가 연접임을 보이면 충분하다
(Descent on Spaces, Lemma
\ref{spaces-descent-lemma-finite-type-descends}를 보라). $f$가 평탄하므로
이는 $f^*\mathcal{F} \cap f^*f_*\mathcal{G}$와 같다. $f$가 평탄하고
준콤팩트이며 준분리이므로
$f^*f_*\mathcal{G} = p_*q^*\mathcal{G}$이다. 여기서
$p, q : Y \times_X Y \to Y$는 사영이다. Cohomology of Spaces, Lemma
\ref{spaces-cohomology-lemma-flat-base-change-cohomology}를 보라. $p$는
절단을 가지므로 결론이 따른다.
\end{proof}

\noindent
Section \ref{spaces-groupoids-section-notation}에서와 같이 $B \to S$라 하자.
$(U, R, s, t, c)$를 대수공간에서의 groupoid라 하고 $B$ 위에 두자.
$U$가 국소 뇌터라고 가정하자. 아래 보조정리에서는
준연접층 $(\mathcal{F}, \alpha)$가 $(U, R, s, t, c)$ 위에 있을 때
{\it 연접}이라는 말을 $\mathcal{F}$가 연접 $\mathcal{O}_U$-가군이라는
뜻으로 쓴다.

\begin{lemma}
\label{spaces-groupoids-lemma-colimit-coherent}
Section \ref{spaces-groupoids-section-notation}에서와 같이 $B \to S$라 하자.
$(U, R, s, t, c)$를 대수공간에서의 groupoid라 하고 $B$ 위에 두자.
다음을 가정하자.
\begin{enumerate}
\item $U$, $R$은 뇌터이다.
\item $s, t$는 평탄하고 준콤팩트이며 준분리이다.
\end{enumerate}
그러면 모든 준연접 가군 $(\mathcal{F}, \alpha)$가
$(U, R, s, t, c)$ 위에 있을 때 연접 가군들의 여과 여극한이다.
\end{lemma}

\begin{proof}
국소 뇌터 대수공간 위 연접 가군에 대한 Cohomology of Spaces,
Lemma \ref{spaces-cohomology-lemma-coherent-Noetherian}의 특성화를
더 언급하지 않고 쓰겠다. $\mathcal{F} = \colim \mathcal{H}_i$를
연접 부분가군 $\mathcal{H}_i \subset \mathcal{F}$들의 여과 여극한으로
쓸 수 있다. Cohomology of Spaces, Lemma
\ref{spaces-cohomology-lemma-directed-colimit-coherent}를 보라. 준연접층
$\mathcal{H}$가 $U$ 위에 주어지면, Lemma
\ref{spaces-groupoids-lemma-construct-quasi-coherent}의 준연접층
$(s_*t^*\mathcal{H}, \alpha)$를 $(U, R, s, t, c)$ 위에 있는 것으로
나타내겠다. 수반 사상
$(\mathcal{F}, \beta) \to (s_*t^*\mathcal{F}, \alpha)$를
$\QCoh(U, R, s, t, c)$에서 생각하자. Remark
\ref{spaces-groupoids-remark-adjunction-map}을 보라. 다음과 같이 둔다.
$$
(\mathcal{F}_i, \beta_i) =
(\mathcal{F}, \beta)
\times_{(s_*t^*\mathcal{F}, \alpha)}
(s_*t^*\mathcal{H}_i, \alpha)
$$
이는 $\QCoh(U, R, s, t, c)$에서의 정의이다. $U$로의 제한은
Lemma \ref{spaces-groupoids-lemma-abelian}의 증명에 의해
$\QCoh(U, R, s, t, c)$ 위의 완전함자이므로 올곱 도식
$$
\xymatrix{
\mathcal{F} \ar[r] & s_*t^*\mathcal{F} \\
\mathcal{F}_i \ar[r] \ar[u] & s_*t^*\mathcal{H}_i \ar[u]
}
$$
을 얻는다. 달리 말하면
$\mathcal{F}_i = \mathcal{F} \cap s_*t^*\mathcal{H}_i$이다.
Remark \ref{spaces-groupoids-remark-adjunction-map}의 수반 사상 설명에 의해 이 도식은
도식
$$
\xymatrix{
\mathcal{F} \ar[r] & s_*s^*\mathcal{F} \\
\mathcal{F}_i \ar[r] \ar[u] & s_*t^*\mathcal{H}_i \ar[u]
}
$$
과 동형이다. 여기서 오른쪽 수직 화살표는 $s_*$를 사상
$$
t^*\mathcal{H}_i \to t^*\mathcal{F} \xrightarrow{\beta} s^*\mathcal{F}
$$
에 적용한 결과이다. $t$가 평탄 사상이므로 이 화살표는
단사이다. 따라서 Lemma \ref{spaces-groupoids-lemma-push-pull}에 의해
$\mathcal{F}_i$는 연접이다. 끝으로 $s$가 준콤팩트이고 준분리이므로
$s_*$는 여극한과 가환한다(Cohomology of Schemes, Lemma
\ref{coherent-lemma-colimit-cohomology}를 보라). 따라서
$s_*t^*\mathcal{F} = \colim s_*t^*\mathcal{H}_i$이고, 원하는 대로
$(\mathcal{F}, \beta) = \colim (\mathcal{F}_i, \beta_i)$이다.
\end{proof}








\section{준연접층의 결정}
\label{spaces-groupoids-section-crystals}

\noindent
$(I, \Phi, j)$를 집합 $I$와 전관계 $j : \Phi \to I \times I$로
이루어진 쌍이라 하자. 각 $i \in I$에 대해 스킴 $X_i$가, 각
$\phi \in \Phi$에 대해 스킴의 사상 $f_\phi : X_{i'} \to X_i$가
주어졌다고 하자. 여기서 $j(\phi) = (i, i')$이다. 다음과 같이 둔다.
$X = (\{X_i\}_{i \in I}, \{f_\phi\}_{\phi \in \Phi})$.
{\it $X$ 위 준연접 가군의 결정}이란 다음 규칙을 말한다. 이 규칙은 각
% P09-ERR-0890
$i \in \Ob(\mathcal{I})$에 준연접층 $\mathcal{F}_i$를 대응시키고,
이 층은 $X_i$ 위에 있다. 또한 각 $\phi \in \Phi$ 가운데
$j(\phi) = (i, i')$인 것에 준연접층의 동형
$$
\alpha_\phi : f_\phi^*\mathcal{F}_i \longrightarrow \mathcal{F}_{i'}
$$
을 대응시키고, 이 동형은 $X_{i'}$ 위에 있다. 이러한 준연접 가군의 결정들은 가법 범주
$\textit{CQC}(X)$를 이룬다.\footnote{삼중항들의 집합
$\phi, \phi', \phi'' \in \Phi$를 따로 정하여
$j(\phi) = (i, i')$, $j(\phi') = (i', i'')$,
$j(\phi'') = (i, i'')$ 및
$f_{\phi''} = f_\phi \circ f_{\phi'}$를 만족하게 하고, 이 삼중항들에
대하여 $\alpha_{\phi'} \circ f_{\phi'}^*\alpha_\phi = \alpha_{\phi''}$를
요구할 수도 있다. 그러면 가법 부분범주가 정의된다. 예를 들어
$(I, \Phi)$라는 자료는 지표 범주의 대상과 화살표의 집합일 수 있고,
$X$는 이 지표 범주 위 스킴의 도식일 수 있다. Lemma
\ref{spaces-groupoids-lemma-crystals-in-quasi-coherent-modules}의 결론은 그 부분범주에서도
곧바로 대응하는 결론을 준다.} 이 범주는 여극한을 가진다(증명은
Lemma \ref{spaces-groupoids-lemma-colimits}의 증명과 같다). 모든 사상 $f_\phi$가
평탄하면 $\textit{CQC}(X)$는 아벨 범주이다(증명은 Lemma
\ref{spaces-groupoids-lemma-abelian}의 증명과 같다). $\kappa$를 기수라 하자. 준연접
가군의 결정 $\mathcal{F}$가 $X$ 위에 있고
{\it $\kappa$-생성}이라는 것은
각 $\mathcal{F}_i$가 $\kappa$-생성이라는 뜻이다(Properties,
Definition \ref{properties-definition-kappa-generated}을 보라).

\begin{lemma}
\label{spaces-groupoids-lemma-crystals-in-quasi-coherent-modules}
위의 상황에서 모든 사상 $f_\phi$가 평탄하면, 어떤 기수 $\kappa$가
존재하여 모든 대상
$(\{\mathcal{F}_i\}_{i \in I}, \{\alpha_\phi\}_{\phi \in \Phi})$가
$\textit{CQC}(X)$ 안에서
자신의 $\kappa$-생성 부분가군들의 유향 여극한이 된다.
\end{lemma}

\begin{proof}
이 보조정리와 증명에서
$(\{\mathcal{F}_i\}_{i \in I}, \{\alpha_\phi\}_{\phi \in \Phi})$의
{\it 부분가군}이란 준연접 부분가군
$\mathcal{G}_i \subset \mathcal{F}_i$를 모든 $i$에 대해 주어
$\alpha_\phi(f_\phi^*\mathcal{G}_i) = \mathcal{G}_{i'}$가 모든
$\mathcal{F}_{i'}$의 부분층으로서 모든 $\phi \in \Phi$에 대해 성립하게 하는
자료를 뜻한다. $f_\phi$가 평탄하면 당김 $f^*_\phi$는 완전함자, 즉
부분층을 보존하므로 이 정의는 타당하다.
% P09-ERR-0891
이 증명은 Properties, Lemma \ref{properties-lemma-colimit-kappa}의 증명의
한 변형이다. 먼저 그 증명을 읽기를 권한다.

\medskip\noindent
모든 스킴 $X_i$가 아핀인 경우에 보조정리를 증명하면 충분하다고
주장한다. 이를 보기 위해
$$
J = \coprod\nolimits_{i \in I} \{U \subset X_i\text{ affine open}\}
$$
로 두고,
% P09-ERR-0892
\begin{align*}
\Psi = & \coprod\nolimits_{\phi \in \Phi}
\{
(U, V) \mid
U \subset X_i, V \subset X_{i'}\text{ affine open with } f_\phi(U) \subset V
\} \\
&
\amalg \coprod\nolimits_{i \in I}
\{
(U, U') \mid
U, U' \subset X_i\text{ affine open with } U \subset U'
\}
\end{align*}
에 자명한 사상 $\Psi \to J \times J$를 갖추자. 그러면
$(\mathcal{F}, \alpha)$는 준연접층의 결정
$(\{\mathcal{H}_j\}_{j \in J}, \{\beta_\psi\}_{\psi \in \Psi})$를
유도하고, 이 결정은 $Y = (J, \Psi)$ 위에 있다. 이는
$\mathcal{H}_{(i, U)} = \mathcal{F}_i|_U$로 두되
$(i, U) \in J$에 대해 그렇게 하고, $\beta_\psi$를
$\psi \in \Psi$에 대해 다음과 같이 두어 얻는다.
% P09-ERR-0893
$\alpha_\phi$를 $U$에 제한한 것으로 두되 이는
$\psi = (\phi, U, V)$일 때이고,
$\text{id} : (\mathcal{F}_i|_{U'})|_U \to \mathcal{F}_i|_U$로 두되 이는
$\psi = (i, U, U')$일 때이다.
더 나아가
$(\{\mathcal{H}_j\}_{j \in J}, \{\beta_\psi\}_{\psi \in \Psi})$의
부분가군들은 $1$-대-$1$로
$(\{\mathcal{F}_i\}_{i \in I}, \{\alpha_\phi\}_{\phi \in \Phi})$의
부분가군들과 일대일로 대응한다. 증명은 생략한다(힌트: Sheaves,
Section \ref{sheaves-section-bases}를 사용하라). 또한 $\kappa$가 $Y$에
대해 성립하면 같은 $\kappa$가 $X$에 대해서도 성립함이 분명하다
($\kappa$-생성 가군의 정의에 의한다).
% P09-ERR-0894
따라서 $Y$ 위 준연접층의 결정에 대하여 보조정리를 증명하면 충분하다.

\medskip\noindent
모든 스킴 $X_i$가 아핀이라고 가정하자. $\kappa$를 $I$ 또는 $\Phi$의
크기보다 큰 무한 기수라 하자. 대상
$(\{\mathcal{F}_i\}_{i \in I}, \{\alpha_\phi\}_{\phi \in \Phi})$를
$\textit{CQC}(X)$에서 잡자. 각 $i$에 대하여 $X_i = \Spec(A_i)$ 및
$M_i = \Gamma(X_i, \mathcal{F}_i)$로 쓰자. 모든 $\phi \in \Phi$ 가운데
$j(\phi) = (i, i')$인 것에 대하여 $\alpha_\phi$는 $A_{i'}$-가군 동형
$$
\alpha_\phi : M_i \otimes_{A_i} A_{i'} \longrightarrow M_{i'}
$$
으로 옮겨진다. 선택 공리를 사용하여 규칙
% P09-ERR-0895
$$
(\phi, m) \longmapsto S(\phi, m')
$$
을 택하자. 정의역은 쌍 $(\phi, m')$의 모음이다. 여기서
$\phi \in \Phi$이고 $j(\phi) = (i, i')$이며 $m' \in M_{i'}$이다.
값은 유한 부분집합
$S(\phi, m') \subset M_i$이며
$$
m' = \alpha_\phi\left(\sum\nolimits_{m \in S(\phi, m')} m \otimes a'_m\right)
$$
이 어떤 $a'_m \in A_{i'}$에 대해 성립한다.

\medskip\noindent
이 선택들을 한 뒤, 임의의 $\mathcal{F}_i$의 임의의 $X_i$ 위 절단이
$\kappa$-생성 부분가군에 속한다고 주장한다. 이를 보기 위해 부분집합들의
모음 $\mathcal{S} = \{S_i\}_{i \in I}$가 주어지고 각
$S_i \subset M_i$의 크기가 $\kappa$ 이하라고 하자. 새 모음
$\mathcal{S}' = \{S'_i\}_{i \in I}$을 다음과 같이 정의한다.
$$
S'_i = S_i \cup
\bigcup\nolimits_{(\phi, m'),\ j(\phi) = (i, i'),\ m' \in S_{i'}} S(\phi, m')
$$
각 $S'_i$의 크기는 여전히 $\kappa$ 이하이다.
$\mathcal{S}^{(0)} = \mathcal{S}$,
$\mathcal{S}^{(1)} = \mathcal{S}'$로
두고 귀납적으로 $\mathcal{S}^{(n + 1)} = (\mathcal{S}^{(n)})'$로 둔다.
그런 다음 $S_i^{(\infty)} = \bigcup_{n \geq 0} S_i^{(n)}$ 및
$\mathcal{S}^{(\infty)} = \{S_i^{(\infty)}\}_{i \in I}$로 둔다.
구성상 모든 $\phi \in \Phi$ 가운데 $j(\phi) = (i, i')$인 것과 모든
$m' \in S^{(\infty)}_{i'}$에 대하여 $m'$을
$\alpha_\phi(m \otimes 1)$ 꼴의 상들의 유한 선형결합으로 쓸 수 있으며,
여기서 $m \in S_i^{(\infty)}$이다. 따라서 $N_i$를
$A_i$-부분가군으로 두되, 이는 $M_i$ 안에서 $S_i^{(\infty)}$가 생성한
것으로 한다. 그러면 대응하는
준연접 부분가군
% P09-ERR-0896
$\widetilde{N_i} \subset \mathcal{F}_i$들은 $\kappa$-생성 부분가군을
이룬다. 이로써 증명이 끝난다.
\end{proof}

\begin{lemma}
\label{spaces-groupoids-lemma-set-generators}
Section \ref{spaces-groupoids-section-notation}에서와 같이 $B \to S$라 하자.
$(U, R, s, t, c)$를 대수공간에서의 groupoid라 하고 $B$ 위에 두자.
$s$, $t$가 평탄하면 어떤 집합 $T$와
대상들의 족 $(\mathcal{F}_t, \alpha_t)_{t \in T}$가
$\QCoh(U, R, s, t, c)$ 안에 존재하여, 모든 대상
$(\mathcal{F}, \alpha)$가 이 대상들 가운데 하나와 동형인 자신의
부분가군들, 즉 $(\mathcal{F}_t, \alpha_t)$ 가운데 하나와 동형인
부분가군들의 유향 여극한이 된다.
\end{lemma}

\begin{proof}
이 보조정리는 스킴에서의 groupoid를 다루는 Groupoids, Lemma
\ref{groupoids-lemma-colimit-kappa}의 일반화이다. 정확히 같은 논증을
쓸 수는 없으므로, 위에서 전개한 ``준연접층의 결정''에 관한 내용을
사용한다.

\medskip\noindent
스킴 $W$와 전사 에탈 사상 $W \to U$를 택하자. 스킴 $V$와 전사 에탈 사상
$V \to W \times_{U, s} R$를 택하자. 스킴 $V'$과 전사 에탈 사상
$V' \to R \times_{t, U} W$를 택하자. 스킴들의 모음
$$
I = \{W, W \times_U W, V, V', V \times_R V'\}
$$
과 스킴 사상들의 집합
$$
\Phi = \{\text{pr}_i : W \times_U W \to W, V \to W, V' \to W,
V \times_R V' \to V, V \times_R V' \to V'\}
$$
을 생각하자. $X = (I, \Phi)$로 둔다. 준연접층의 결정의 범주
$\textit{CQC}(X)$를 정의했음을 상기하자. 이 범주의 결정들은 $X$ 위에 있다.
다음 함자가 있다.
$$
\QCoh(U, R, s, t, c) \longrightarrow \textit{CQC}(X)
$$
이 함자는 $(\mathcal{F}, \alpha)$에 다음 자료를 대응시킨다. 층
$\mathcal{F}|_W$를
$W$ 위에, 층 $\mathcal{F}|_{W \times_U W}$를 $W \times_U W$ 위에,
$\mathcal{F}$를 $V \to W \times_{U, s} R \to W \to U$를 통해 당겨온
층을 $V$ 위에, $\mathcal{F}$를
$V' \to R \times_{t, U} W \to W \to U$를 통해 당겨온 층을 $V'$ 위에,
끝으로 $\mathcal{F}$를
$V \times_R V' \to V \to W \times_{U, s} R \to W \to U$를 통해
당겨온 층을 $V \times_R V'$ 위에 둔다. 비교 사상
$\{\alpha_\phi\}_{\phi \in \Phi}$에는 당김의 결합법칙에서 오는
자명한 것들을 쓰되, 사상
$\phi = \text{pr}_{V'} : V \times_R V' \to V'$에는
$\alpha : t^*\mathcal{F} \to s^*\mathcal{F}$를 $V \times_R V'$로
당겨온 것을 쓴다. 이는 다음 가환 도식 때문에 타당하다.
$$
\xymatrix{
& V \times_R V' \ar[ld] \ar[rd] \\
V \ar[rd] \ar[dd] & & V' \ar[ld] \ar[dd] \\
& R \ar@<-1ex>[dd]_s \ar@<1ex>[dd]^t \\
W \ar[rd] & & W \ar[ld] \\
& U
}
$$
위에 표시한 함자는 범주의 동치가 아니다. 그러나 $W \to U$가 전사
에탈이므로 이 함자는 충실하다.\footnote{실제로 이 함자는 완전충실하지만,
여기서는 필요하지 않다.} 위 도식의 모든 사상이 평탄하므로 이 함자는
아벨 범주 사이의 완전함자이다. 또한
% P09-ERR-0897
$(\mathcal{F}, \alpha)$가 주어지고 그 상이
$(\{\mathcal{F}_i\}_{i \in I}, \{\alpha_\phi\}_{\phi \in \Phi})$이면
$1$-대-$1$ 대응이 존재하며, 이는
$(\mathcal{F}, \alpha)$의 준연접 부분가군들과
$(\{\mathcal{F}_i\}_{i \in I}, \{\alpha_\phi\}_{\phi \in \Phi})$의
부분가군들 사이의 대응이라고 주장한다. 실제로
$(\{\mathcal{F}_i\}_{i \in I}, \{\alpha_\phi\}_{\phi \in \Phi})$의
부분가군이 주어졌을 때, $W$ 위 부분가군과 사영 사상
$W \times_U W \to W$의
양립성은 이 부분가군이 $\mathcal{F}$의 준연접 부분가군에서 옴을
보장한다(Properties of Spaces, Proposition
\ref{spaces-properties-proposition-quasi-coherent}에 의한다). 또한
$\alpha_{\text{pr}_{V'}}$와의 양립성은 이 부분층이 $\alpha$와 양립함을
보장한다(세부사항은 생략한다).

\medskip\noindent
기수 $\kappa$를 Lemma \ref{spaces-groupoids-lemma-crystals-in-quasi-coherent-modules}에서처럼
계 $X = (I, \Phi)$에 대해 택하자. Properties, Lemma
\ref{properties-lemma-set-of-iso-classes}에서 $\kappa$-생성인
준연접층의 결정 가운데 $X$ 위에 있는 것들의 동형류들이 집합을 이룸이
분명하다. 따라서 결론이
분명하다.
\end{proof}
\section{Groupoid와 군 공간}
\label{spaces-groupoids-section-groupoids-group-spaces}

\noindent
Groupoids, Section \ref{groupoids-section-groupoids-group-schemes}와
비교하라.

\begin{lemma}
\label{spaces-groupoids-lemma-groupoid-from-action}
Section \ref{spaces-groupoids-section-notation}에서와 같은 $B \to S$를 잡자.
군 대수공간 $(G, m)$을 $B$ 위에 잡되 그 항등원을 $e_G$, 역원을
$i_G$라 하자. 대수공간 $X$를 $B$ 위에 잡고, 사상
$a : G \times_B X \to X$를 $G$의 $X$에 대한 $B$ 위 작용이라
하자. 그러면 대수공간에서의 groupoid $(U, R, s, t, c, e, i)$를
$B$ 위에서 다음과 같이 얻는다.
\begin{enumerate}
\item $U = X$, $R = G \times_B X$로 둔다.
\item $s : R \to U$를 $(g, x) \mapsto x$로 둔다.
\item $t : R \to U$를 $(g, x) \mapsto a(g, x)$로 둔다.
\item $c : R \times_{s, U, t} R \to R$를
$((g, x), (g', x')) \mapsto (m(g, g'), x')$로 둔다.
\item $e : U \to R$를 $x \mapsto (e_G(x), x)$로 둔다.
\item $i : R \to R$를 $(g, x) \mapsto (i_G(g), a(g, x))$로 둔다.
\end{enumerate}
\end{lemma}

\begin{proof}
생략한다. 힌트: 집합 수준에서 이것이 성립함을 보이면 충분하다.
이를 위해 보조정리 앞에서 $g$를 $v$에서 $a(g, v)$로 가는
화살표로 설명한 내용을 사용하라.
\end{proof}

\begin{lemma}
\label{spaces-groupoids-lemma-action-groupoid-modules}
Section \ref{spaces-groupoids-section-notation}에서와 같은 $B \to S$를 잡자.
군 대수공간 $(G, m)$을 $B$ 위에 잡자.
대수공간 $X$를 $B$ 위에 잡고, 사상
$a : G \times_B X \to X$를 $G$의 $X$에 대한 $B$ 위 작용이라
하자. Lemma
\ref{spaces-groupoids-lemma-groupoid-from-action}에서 구성한 대수공간에서의
groupoid를 $(U, R, s, t, c)$라 하자.
규칙
$(\mathcal{F}, \alpha) \mapsto (\mathcal{F}, \alpha)$는
$G$-동변 $\mathcal{O}_X$-가군의 범주와
$(U, R, s, t, c)$ 위 준연접 가군의 범주 사이의 범주 동치를 정의한다.
\end{lemma}

\begin{proof}
이 주장은 의미가 있다. 실제로 $t = a$이고 $s = \text{pr}_1$이며,
이들은 사상 $R = G \times_B X \to X$이다. Definitions
\ref{spaces-groupoids-definition-equivariant-module}와
\ref{spaces-groupoids-definition-groupoid-module}를 보라.
Lemma \ref{spaces-groupoids-lemma-groupoid-from-action}의 번역을 사용하면 두 정의의
가환성 조건이 정확히 서로 일치한다.
\end{proof}




\section{안정자 군 대수공간}
\label{spaces-groupoids-section-stabilizer}

\noindent
Groupoids, Section \ref{groupoids-section-stabilizer}와 비교하라.
대수공간에서의 groupoid가 주어지면 다음과 같이 군 대수공간을 얻는다.

\begin{lemma}
\label{spaces-groupoids-lemma-groupoid-stabilizer}
Section \ref{spaces-groupoids-section-notation}에서와 같은 $B \to S$를 잡자.
대수공간에서의 groupoid $(U, R, s, t, c)$를 $B$ 위에 잡자.
다음 데카르트 정사각형으로 정의되는 대수공간 $G$
$$
\xymatrix{
G \ar[r] \ar[d] & R \ar[d]^{j = (t, s)} \\
U \ar[r]^-{\Delta} & U \times_B U
}
$$
는 $U$ 위 군 대수공간이며, 그 합성법칙 $m$은 합성법칙 $c$에 의해
유도된다.
\end{lemma}

\begin{proof}
이는 groupoid 범주에서 임의의 대상의 자기 사상들의 집합이 군을
이루기 때문이다.
\end{proof}

\noindent
$\Delta$가 단사사상이므로 $G = j^{-1}(\Delta_{U/B})$는 $R$의
부분층이다. 이렇게 생각하면 구조 사상
$G = j^{-1}(\Delta_{U/B}) \to U$는 $s$ 또는 $t$ 중 어느 것에
의해서도 유도되며(둘은 같다), $m$은 $c$에 의해 유도된다.

\begin{definition}
\label{spaces-groupoids-definition-stabilizer-groupoid}
Section \ref{spaces-groupoids-section-notation}에서와 같은 $B \to S$를 잡자.
대수공간에서의 groupoid $(U, R, s, t, c)$를 $B$ 위에 잡자.
군 대수공간 $j^{-1}(\Delta_{U/B}) \to U$를
{\it 대수공간에서의 groupoid $(U, R, s, t, c)$의 안정자}라 한다.
\end{definition}

\noindent
문헌에서는 안정자 군 대수공간을 흔히 $S$로 나타낸다(아마도
stabilizer라는 단어가 s로 시작하기 때문일 것이다). 여기서는
이미 기저 스킴을 $S$로 나타냈으므로 그렇게 할 수 없다.

\begin{lemma}
\label{spaces-groupoids-lemma-groupoid-action-stabilizer}
Section \ref{spaces-groupoids-section-notation}에서와 같은 $B \to S$를 잡자.
대수공간에서의 groupoid $(U, R, s, t, c)$를 $B$ 위에 잡고 그
안정자를 $G/U$라 하자. $R_t/U$는 대수공간 $R$을 $U$ 위에서
사상 $t : R \to U$를 통하여 본 것이라고 나타내자.
표준적인 왼쪽 작용
$$
a : G \times_U R_t \longrightarrow R_t
$$
이 존재하며, 이는 합성법칙 $c$에 의해 유도된다.
\end{lemma}

\begin{proof}
$T/B$ 위의 점으로 표현하면 $a(g, r) = c(g, r)$로 정의한다.
\end{proof}







\section{Groupoid의 제한}
\label{spaces-groupoids-section-restrict-groupoid}

\noindent
표기법은 Groupoids, Section
\ref{groupoids-section-restrict-groupoid}을 참조하라.

\begin{lemma}
\label{spaces-groupoids-lemma-restrict-groupoid}
Section \ref{spaces-groupoids-section-notation}에서와 같은 $B \to S$를 잡자.
대수공간에서의 groupoid $(U, R, s, t, c)$를 $B$ 위에 잡자.
대수공간의 사상 $g : U' \to U$를 잡자.
다음 도식을 생각하자.
$$
\xymatrix{
R' \ar[d] \ar[r] \ar@/_3pc/[dd]_{t'} \ar@/^1pc/[rr]^{s'}&
R \times_{s, U} U' \ar[r] \ar[d] &
U' \ar[d]^g \\
U' \times_{U, t} R \ar[d] \ar[r] &
R \ar[r]^s \ar[d]_t &
U \\
U' \ar[r]^g &
U
}
$$
여기서 모든 정사각형은 섬유곱 정사각형이다. 그러면 표준적인
합성법칙 $c' : R' \times_{s', U', t'} R' \to R'$가 존재하여
$(U', R', s', t', c')$는 $B$ 위 대수공간에서의 groupoid가 되고,
$U' \to U$, $R' \to R$은 사상
$(U', R', s', t', c') \to (U, R, s, t, c)$를 $B$ 위
대수공간에서의 groupoid 사이에 정의한다. 더욱이 임의의 스킴
$T$를 $B$ 위에 잡으면 groupoid의 함자
$$
(U'(T), R'(T), s', t', c') \to (U(T), R(T), s, t, c)
$$
는 $(U(T), R(T), s, t, c)$를 사상 $U'(T) \to U(T)$를 통하여
제한한 것이다(Groupoids, Section
\ref{groupoids-section-restrict-groupoid} 참조).
\end{lemma}

\begin{proof}
생략한다.
\end{proof}

\begin{definition}
\label{spaces-groupoids-definition-restrict-groupoid}
Section \ref{spaces-groupoids-section-notation}에서와 같은 $B \to S$를 잡자.
대수공간에서의 groupoid $(U, R, s, t, c)$를 $B$ 위에 잡자.
사상 $g : U' \to U$를 대수공간의 사상으로서 $B$ 위에 잡자.
Lemma \ref{spaces-groupoids-lemma-restrict-groupoid}에서 구성한 대수공간에서의
groupoid의 사상
$(U', R', s', t', c') \to (U, R, s, t, c)$를
{\it $(U, R, s, t, c)$의 $U'$로의 제한}이라 한다.
이 경우 때때로 $R' = R|_{U'}$라는 표기법을 사용한다.
\end{definition}

\begin{lemma}
\label{spaces-groupoids-lemma-restrict-groupoid-relation}
Definitions \ref{spaces-groupoids-definition-restrict-groupoid}와
\ref{spaces-groupoids-definition-restrict-relation}에서 정의한 groupoid와
(준)동치관계의 제한 개념은 Lemmas
\ref{spaces-groupoids-lemma-groupoid-pre-equivalence}와
\ref{spaces-groupoids-lemma-equivalence-groupoid}의 구성을 통하여 일치한다.
\end{lemma}

\begin{proof}
여기서 말하는 것은 Lemma \ref{spaces-groupoids-lemma-restrict-groupoid}의 $R'$가
또한
$$
R' = (U' \times_B U')\times_{U \times_B U} R
\longrightarrow
U' \times_B U'
$$
와 같다는 것이다. 실제로 이것이 그 보조정리를 진술하는 더 명료한
방식이었을 수도 있다.
\end{proof}





\section{불변 부분공간}
\label{spaces-groupoids-section-invariant}

\noindent
이 절에서는 불변 부분공간의 개념을 간단히 논한다.

\begin{definition}
\label{spaces-groupoids-definition-invariant-open}
Section \ref{spaces-groupoids-section-notation}에서와 같은 $B \to S$를 잡자.
대수공간에서의 groupoid $(U, R, s, t, c)$를 기저 $B$ 위에 잡자.
\begin{enumerate}
\item 열린 부분공간 $W \subset U$가 {\it $R$-불변}이라는 것은
$t(s^{-1}(W)) \subset W$를 만족한다는 뜻이다.
\item 국소 닫힌 부분공간 $Z \subset U$가 $R$의 국소 닫힌
부분공간으로서 $t^{-1}(Z) = s^{-1}(Z)$를 만족하면
{\it $R$-불변}이라고 한다.
\item 대수공간의 단사사상 $T \to U$가 $R$ 위 대수공간으로서
$T \times_{U, t} R = R \times_{s, U} T$를 만족하면
{\it $R$-불변}이라고 한다.
\end{enumerate}
\end{definition}

\noindent
열린 부분공간 $W \subset U$의 $R$-불변성은
$s^{-1}(W) = t^{-1}(W)$를 요구하는 것과도 동치이다.
$W \subset U$가 $R$-불변이면 $R$의 $W$로의 제한은 단지
$R_W = s^{-1}(W) = t^{-1}(W)$이다. 마찬가지로 $Z \subset U$가
$R$-불변 국소 닫힌 부분공간이면 $R$의 $Z$로의 제한은 단지
$R_Z = s^{-1}(Z) = t^{-1}(Z)$이다.

\begin{lemma}
\label{spaces-groupoids-lemma-constructing-invariant-opens}
Section \ref{spaces-groupoids-section-notation}에서와 같은 $B \to S$를 잡자.
대수공간에서의 groupoid $(U, R, s, t, c)$를 $B$ 위에 잡자.
\begin{enumerate}
\item $s$와 $t$가 열린 사상이면, 모든 열린집합 $W \subset U$에
대하여 열린집합 $s(t^{-1}(W))$는 $R$-불변이다.
\item $s$와 $t$가 열린 준콤팩트 사상이면, $U$는 $R$-불변
준콤팩트 열린 부분공간들로 이루어진 열린 덮개를 갖는다.
\end{enumerate}
\end{lemma}

\begin{proof}
$s$와 $t$가 열린 사상이고 $W \subset U$가 열린집합이라고 하자.
$s$가 열린 사상이므로 $W' = s(t^{-1}(W))$는 $U$의 열린
부분공간이다. 이제 함자적 관점으로 이것이 $R$-불변인 $U$의
열린 부분집합임을 확인하기는 매우 쉽지만, 도식으로 직접 논증하는
것이 유익하다고 생각하므로 그렇게 하겠다.
$t^{-1}(W')$는 사상
$$
A := t^{-1}(W) \times_{s|_{t^{-1}(W)}, U, t} R
\xrightarrow{\text{pr}_1} R
$$
의 상이고, $s^{-1}(W')$는 사상
$$
B := R \times_{s, U, s|_{t^{-1}(W)}} t^{-1}(W)
\xrightarrow{\text{pr}_0} R.
$$
의 상임에 유의하라. 위 화살표 왼쪽의 대수공간 $A$, $B$는 각각
$R \times_{s, U, t} R$와 $R \times_{s, U, s} R$의 열린
부분공간이다. Lemma \ref{spaces-groupoids-lemma-diagram}에 의해 도식
$$
\xymatrix{
R \times_{s, U, t} R \ar[rd]_{\text{pr}_1} \ar[rr]_{(\text{pr}_1, c)} & &
R \times_{s, U, s} R \ar[ld]^{\text{pr}_0} \\
& R &
}
$$
은 가환하고 수평 화살표는 동형사상이다. 더욱이
$(\text{pr}_1, c)(A) = B$임이 명백하다. 따라서
$s^{-1}(W') = t^{-1}(W')$이고 $W'$는 $R$-불변이다.
이로써 (1)을 증명했다.

\medskip\noindent
이제 $s$, $t$가 모두 열린 준콤팩트 사상이라고 하자.
$W \subset U$가 준콤팩트 열린집합이면
$W' = s(t^{-1}(W))$도 준콤팩트 열린집합이고, 위 논의에 의해
불변이다. $W$가 $U$ 위 에탈 아핀들의 상들을 모두 달리하게 하면
(2)를 얻는다.
\end{proof}





\section{몫 층}
\label{spaces-groupoids-section-quotient-sheaves}

\noindent
$S$를 스킴이라 하고 $B$를 $S$ 위 대수공간이라 하자.
$j : R \to U \times_B U$를 $B$ 위 준관계라 하자.
각 스킴 $S'$를 $S$ 위에 잡자. 동치관계
$\sim_{S'}$는 사상 $j(S') : R(S') \to U(S') \times U(S')$의
상이 생성하는 것으로 취할 수 있다. 따라서 준층
\begin{equation}
\label{spaces-groupoids-equation-quotient-presheaf}
\begin{matrix}
(\Sch/S)^{opp}_{fppf} &
\longrightarrow &
\textit{Sets}, \\
S' &
\longmapsto  &
U(S')/\sim_{S'}
\end{matrix}
\end{equation}
을 얻는다. $j$는 $B$ 위 대수공간의 사상이고
$U \times_B U$로 가므로, 준층
(\ref{spaces-groupoids-equation-quotient-presheaf})에서 $B$로 가는 표준적인 준층의
변환이 있음에 유의하라.

\begin{definition}
\label{spaces-groupoids-definition-quotient-sheaf}
위와 같은 $B \to S$와 준관계 $j : R \to U \times_B U$를 잡자.
이 상황에서 {\it 몫 층 $U/R$}은 $j$에 연관되며,
$(\Sch/S)_{fppf}$ 위 준층
(\ref{spaces-groupoids-equation-quotient-presheaf})의 층화이다.
$j : R \to U \times_B U$가 Lemma
\ref{spaces-groupoids-lemma-groupoid-from-action}에서와 같이 군 대수공간
$G$의 $B$ 위에서의 $U$에 대한 작용에서 오면, 몫 층을
$U/G$라 나타낸다.
\end{definition}

\noindent
이는 정확히 도식
$$
\xymatrix{
R \ar@<1ex>[r] \ar@<-1ex>[r] &
U \ar[r] &
U/R
}
$$
이 $(\Sch/S)_{fppf}$ 위 집합의 층의 범주에서 쌍대등화자 도식이라는
뜻이다. 다시 층의 표준적인 사상 $U/R \to B$가 있다. 이는
$j$가 $B$ 위 대수공간의 사상으로서 $U \times_B U$로 가기
때문이다.

\begin{remark}
\label{spaces-groupoids-remark-quotient-variant}
위 구성의 한 변형은 함자
$$
\begin{matrix}
(\textit{Spaces}/B)^{opp}_{fppf} &
\longrightarrow &
\textit{Sets}, \\
X &
\longmapsto  &
U(X)/\sim_X
\end{matrix}
$$
를 층화하는 것이다. 여기서 이제
$\sim_X \subset U(X) \times U(X)$는
$j : R(X) \to U(X) \times U(X)$의 상이 생성하는 동치관계이다.
물론 여기서 $U(X) = \Mor_B(X, U)$이고
$R(X) = \Mor_B(X, R)$이다. 실제로 (Topologies of Spaces의 미래
참조를 여기에 삽입)의 동일시를 통하면 결과는 같았을 것이다.
\end{remark}

\begin{definition}
\label{spaces-groupoids-definition-representable-quotient}
Definition \ref{spaces-groupoids-definition-quotient-sheaf}의 상황을 가정하자.
준관계 $j$에 대하여 층 $U/R$가 대수공간이면
{\it 대수공간으로 표현되는 몫}을 갖는다고 한다.
준관계 $j$에 대하여 층 $U/R$가 스킴으로 표현되면
{\it 표현 가능한 몫}을 갖는다고 한다.
대수공간에서의 groupoid $(U, R, s, t, c)$가 $B$ 위에 있을 때
몫 $U/R$가 $j = (t, s)$에 대해 표현 가능하면(각각 대수공간이면)
{\it 표현 가능한 몫}을 갖는다고(각각
{\it 대수공간으로 표현되는 몫}을 갖는다고) 한다.
\end{definition}

\noindent
몫 $U/R$가 $M$으로 표현되면($S$ 위의 스킴이나 대수공간 어느
경우든), 위에서 본 것처럼 표준적인 구조 사상 $M \to B$가 함께
주어진다.

\medskip\noindent
다음 보조정리는 몫을 표현하는 $M$을 특징짓는다.
예를 들어 $U \to M$이 평탄하고 유한 표시이며 전사이고
$R \cong U \times_M U$이면 이를 적용할 수 있다.

\begin{lemma}
\label{spaces-groupoids-lemma-criterion-quotient-representable}
Definition \ref{spaces-groupoids-definition-quotient-sheaf}의 상황을 가정하자.
$M$이라는 대수공간이 $S$ 위에 있고 다음을 만족하는 사상
$U \to M$이
존재한다고 가정하자.
\begin{enumerate}
\item 사상 $U \to M$은 $s, t$를 등화한다.
\item 사상 $U \to M$은 층의 전사이다.
\item 유도된 사상 $(t, s) : R \to U \times_M U$은 층의 전사이다.
\end{enumerate}
이 경우 $M$은 몫 층 $U/R$를 표현한다.
\end{lemma}

\begin{proof}
조건 (1)은 $U \to M$이 $U/R$를 통하여 인수분해된다는 뜻이다.
조건 (2)는 층의 사상으로서 $U/R \to M$이 전사라는 뜻이다.
조건 (3)은 층의 사상으로서 $U/R \to M$이 단사라는 뜻이다.
따라서 보조정리를 얻는다.
\end{proof}

\noindent
다음 보조정리는 $j$가 준동치관계라고 요구하지 않으면(가령 단지
준관계라고만 하면) 틀리다.

\begin{lemma}
\label{spaces-groupoids-lemma-quotient-pre-equivalence}
$S$를 스킴이라 하자. $B$를 $S$ 위 대수공간이라 하자.
$j : R \to U \times_B U$를 $B$ 위 준동치관계라 하자.
스킴 $S'$를 $S$ 위에 잡고 $a, b \in U(S')$라 하자. 다음은
동치이다.
\begin{enumerate}
\item $a$와 $b$는 $(U/R)(S')$의 같은 원소로 사상된다.
\item fppf 덮개 $\{f_i : S_i \to S'\}$가 $S'$를 덮고 다음을 만족하는
사상 $r_i : S_i \to R$이 존재한다.
$a \circ f_i = s \circ r_i$이고 $b \circ f_i = t \circ r_i$이다.
\end{enumerate}
달리 말하면, 이 경우 층의 사상
$$
R \longrightarrow U \times_{U/R} U
$$
은 전사이다.
\end{lemma}

\begin{proof}
생략한다. 힌트: 이것이 성립하는 이유는 이 경우 준층
(\ref{spaces-groupoids-equation-quotient-presheaf})가 실제로
$T \mapsto U(T)/j(R(T))$로 주어지기 때문이다.
실제로 $j(R(T)) \subset U(T) \times U(T)$는 동치관계이다.
Definition \ref{spaces-groupoids-definition-equivalence-relation}을 보라.
\end{proof}

\begin{lemma}
\label{spaces-groupoids-lemma-quotient-pre-equivalence-relation-restrict}
$S$를 스킴이라 하자. $B$를 $S$ 위 대수공간이라 하자.
$j : R \to U \times_B U$를 $B$ 위 준동치관계라 하고
$g : U' \to U$를 $B$ 위 대수공간의 사상이라 하자.
$j' : R' \to U' \times_B U'$를 $j$의 $U'$로의 제한이라 하자.
몫 층의 사상
$$
U'/R' \longrightarrow U/R
$$
은 단사이다. $U' \to U$가 층의 사상으로서 전사이면, 예를 들어
$\{g : U' \to U\}$가 fppf 덮개이면(Topologies on Spaces,
Definition \ref{spaces-topologies-definition-fppf-covering} 참조),
$U'/R' \to U/R$는 층의 동형사상이다.
\end{lemma}

\begin{proof}
$\xi, \xi' \in (U'/R')(S')$가 $U/R$의 같은 단면으로 사상되는
단면들이라고 하자. 그러면 fppf 덮개
$\mathcal{S} = \{S_i \to S'\}$를 $S'$에 대해 찾아
$\xi|_{S_i}, \xi'|_{S_i}$가 $a_i, a_i' \in U'(S_i)$로 주어지게 할
수 있다. Lemma \ref{spaces-groupoids-lemma-quotient-pre-equivalence}와 자리의
공리들에 의해 $\mathcal{T}$를 세분한 뒤 사상
$r_i : S_i \to R$이 존재하여
$g \circ a_i = s \circ r_i$, $g \circ a_i' = t \circ r_i$를
만족한다고 가정할 수 있다.
구성상
$R' = R \times_{U \times_S U} (U' \times_S U')$
이므로 $(r_i, (a_i, a_i')) \in R'(S_i)$이고, 이는 $a_i$와
$a_i'$가 $U'/R'$에서 $S_i$ 위의 같은 단면을 정의함을 보인다.
층 조건에 의해 이는 $\xi = \xi'$를 뜻한다.

\medskip\noindent
$U' \to U$가 층의 전사이면 $U'/R' \to U/R$도 전사이다.
마지막으로 $\{g : U' \to U\}$가 fppf 덮개이면 층의 사상
$U' \to U$는 전사이다. Topologies on Spaces, Lemma
\ref{spaces-topologies-lemma-fppf-covering-surjective}를 보라.
\end{proof}

\begin{lemma}
\label{spaces-groupoids-lemma-quotient-groupoid-restrict}
$S$를 스킴이라 하자. $B$를 $S$ 위 대수공간이라 하자.
$(U, R, s, t, c)$를 $B$ 위 대수공간에서의 groupoid라 하자.
$g : U' \to U$를 $B$ 위 대수공간의 사상이라 하자.
$(U', R', s', t', c')$를 $(U, R, s, t, c)$의 $U'$로의 제한이라
하자. 몫 층의 사상
$$
U'/R' \longrightarrow U/R
$$
은 단사이다. 합성
$$
\xymatrix{
U' \times_{g, U, t} R \ar[r]_-{\text{pr}_1} \ar@/^3ex/[rr]^h
& R \ar[r]_s & U
}
$$
이 fppf 층의 전사이면 이 사상은 전단사이다.
예를 들어 $\{h : U' \times_{g, U, t} R \to U\}$가
$fppf$-덮개이거나, $U' \to U$가 층의 전사이거나,
$\{g : U' \to U\}$가 fppf 위상에서 덮개이면 이것이 성립한다.
\end{lemma}

\begin{proof}
단사성은 Lemmas \ref{spaces-groupoids-lemma-groupoid-pre-equivalence}와
\ref{spaces-groupoids-lemma-quotient-pre-equivalence-relation-restrict}를 결합하면
따른다. 전사성을 보이기 위해(Sites, Section
\ref{sites-section-sheaves-injective}에서 층의 전사 사상의
특징짓기를 보라) 다음과 같이 논증한다.
$T$가 스킴이고 $\sigma \in U/R(T)$라고 하자.
덮개 $\{T_i \to T\}$가 존재하여 $\sigma|_{T_i}$가 어떤 원소
$f_i \in U(T_i)$의 상이 된다. 따라서 $\sigma$가
$f \in U(T)$의 상이라고 가정할 수 있다. $h$가 층의 전사라는
가정에 의해 fppf 덮개 $\{\varphi_i : T_i \to T\}$와 사상
$f_i : T_i \to U' \times_{g, U, t} R$를 찾아
$f \circ \varphi_i = h \circ f_i$를 만족하게 할 수 있다.
$f'_i = \text{pr}_0 \circ f_i : T_i \to U'$라 나타내자.
그러면 $f'_i \in U'(T_i)$는 $g \circ f'_i \in U(T_i)$로 사상되고,
$$
g \circ f'_i \sim_{T_i} h \circ f_i = f \circ \varphi_i
$$
이다. 표기법은 (\ref{spaces-groupoids-equation-quotient-presheaf})에서와 같다.
즉 관계를 주는 $R(T_i)$의 원소는 $\text{pr}_1 \circ f_i$이다.
이는 $\sigma$의 $T_i$로의 제한이 원하는 대로
$U'/R'(T_i) \to U/R(T_i)$의 상에 있음을 뜻한다.

\medskip\noindent
$\{h\}$가 fppf 덮개이면 층의 전사를 유도한다. Topologies on
Spaces, Lemma
\ref{spaces-topologies-lemma-fppf-covering-surjective}를 보라.
$U' \to U$가 전사이면 $h$도 전사이다. 실제로 $s$는
단면, 즉 groupoid 스킴의 중립원 $e$를 가진다.
\end{proof}





\section{몫 스택}
\label{spaces-groupoids-section-stacks}

\noindent
이 절과 이어지는 몇 절에서는 위의 Section
\ref{spaces-groupoids-section-quotient-sheaves}와 Groupoids, Section
\ref{groupoids-section-quotient-sheaves}의 한 일반화를 설명한다.
차이는 다음과 같다. 몫 층 대신 몫 스택을 취할 것이다.

\medskip\noindent
스킴 $S$, 대수공간 $B$를 $S$ 위에, 그리고 대수공간에서의
groupoid $(U, R, s, t, c)$를 $B$ 위에 가지고 있다고 하자.
이 자료가 주어지면 함자
\begin{equation}
\label{spaces-groupoids-equation-quotient-stack}
\begin{matrix}
(\Sch/S)_{fppf}^{opp} &
\longrightarrow &
\textit{Groupoids} \\
S' &
\longmapsto &
(U(S'), R(S'), s, t, c)
\end{matrix}
\end{equation}
를 생각한다. Categories, Example
\ref{categories-example-functor-groupoids}에 의해 이 ‘groupoid에서의
준층’은 $(\Sch/S)_{fppf}$ 위 groupoid들로 섬유화된 범주에
대응한다. 이 장에서는 이를
$$
[U/_{\!p}R] \to (\Sch/S)_{fppf}
$$
로 나타낸다. 여기서 아래첨자 $ {}_p$는 몫 스택과 구별하기 위해
붙인 것이다.

\begin{definition}
\label{spaces-groupoids-definition-quotient-stack}
몫 스택. 위와 같은 $B \to S$를 잡자.
\begin{enumerate}
\item $(U, R, s, t, c)$를 $B$ 위 대수공간에서의 groupoid라 하자.
{\it 몫 스택}
$$
p : [U/R] \longrightarrow (\Sch/S)_{fppf}
$$
은 $(U, R, s, t, c)$의 것이며 스택화이다(Stacks, Lemma
\ref{stacks-lemma-stackify-groupoids} 참조). 스택화하는 대상은
groupoid들로 섬유화된 범주 $[U/_{\!p}R]$이며,
$(\Sch/S)_{fppf}$ 위에 있고
(\ref{spaces-groupoids-equation-quotient-stack})에 연관된다.
\item $(G, m)$을 $B$ 위 군 대수공간이라 하자.
$a : G \times_B X \to X$를 $G$의 한 대수공간에 대한 $B$ 위
작용이라 하자. {\it 몫 스택}
$$
p : [X/G] \longrightarrow (\Sch/S)_{fppf}
$$
은 대수공간에서의 groupoid
$(X, G \times_B X, s, t, c)$를 $B$ 위에서 취한 것에 연관된
몫 스택이다. 이는 Lemma \ref{spaces-groupoids-lemma-groupoid-from-action}의
groupoid이다.
\end{enumerate}
\end{definition}

\noindent
따라서 $[U/R]$와 $[X/G]$는 $(\Sch/S)_{fppf}$ 위
groupoid에서의 스택이다. 이 스택들은 뒤에서 매우 중요하므로
상세히 설명하는 것이 타당하다. 대수공간 $X$가 $S$ 위에 주어졌을
때 표기법 $\mathcal{S}_X \to (\Sch/S)_{fppf}$는 층 $X$에
연관된 집합에서의 스택을 나타낸다는 것을 상기하라. Categories,
Lemma \ref{categories-lemma-2-category-fibred-sets}와 Stacks,
Lemma \ref{stacks-lemma-when-stack-in-sets}를 보라.

\begin{lemma}
\label{spaces-groupoids-lemma-quotient-stack-arrows}
Definition \ref{spaces-groupoids-definition-quotient-stack} (1)에서와 같은
$B \to S$와 $(U, R, s, t, c)$를 가정하자.
표준적인 $1$-사상 $\pi : \mathcal{S}_U \to [U/R]$와
$[U/R] \to \mathcal{S}_B$가 존재하며, 이들은
$(\Sch/S)_{fppf}$ 위 groupoid에서의 스택의 사상이다.
합성 $\mathcal{S}_U \to \mathcal{S}_B$는 $1$-사상이며 구조 사상
$U \to B$에 연관된다.
\end{lemma}

\begin{proof}
이 증명에서 준층 (\ref{spaces-groupoids-equation-quotient-stack})에 연관된
groupoid들로 섬유화된 범주를 $[U/_{\!p}R]$로 나타내자.
스택화의 구성에 의해 $1$-사상
$[U/_{\!p}R] \to [U/R]$가 존재한다.
$1$-사상 $\mathcal{S}_U \to [U/R]$는 단순히 합성
$\mathcal{S}_U \to [U/_{\!p}R] \to [U/R]$이다.
여기서 첫 화살표는 스킴 $S'/S$와 사상
$x : S' \to U$를 취하되 이 사상은 $S$ 위에 있고, 대상
$x \in U(S')$를 대응시킨다. 이는 $[U/_{\!p}R]$의 섬유 범주에서
$S'$ 위의 대상이다.

\medskip\noindent
$1$-사상 $[U/R] \to \mathcal{S}_B$를 구성하려면
$1$-사상 $[U/_{\!p}R] \to \mathcal{S}_B$를 구성하면 충분하다.
Stacks, Lemma
\ref{stacks-lemma-stackify-groupoids-universal-property}를 보라.
$S'/S$ 위의 대상에는 사상
$$
U(S') \longrightarrow B(S')
$$
을 사용한다. 이는 구조 사상 $U \to B$에서 온다.
그리고 $a \in R(S')$가 시작점
$s(a) \in U(S')$와 끝점 $t(a) \in U(S')$를 갖는 ‘화살표’라
하자. $s$와 $t$는 $B$ {\it 위의} 사상이므로 두 점은 모두
같은 원소 $\overline{a}$로 사상되며, 이는 $B(S')$의 원소이다.
따라서 화살표
$a \in R(S')$를 $\overline{a}$의 항등사상으로 보낼 수 있다.
(섬유 범주 $(\mathcal{S}_B)_{S'}$에는 항등사상만 있으므로
이것이 옳다.) 이 규칙이 이러한 분할 섬유화 범주에서의 당김과
양립하여 원하는 $1$-사상
$[U/_{\!p}R] \to \mathcal{S}_B$를 정의한다는 확인은 생략한다.

\medskip\noindent
마지막 명제의 확인은 생략한다.
\end{proof}

\begin{lemma}
\label{spaces-groupoids-lemma-quotient-stack-2-arrow}
Lemma \ref{spaces-groupoids-lemma-quotient-stack-arrows}와 같은 가정과 표기법을
사용하자. 표준적인 $2$-사상
$\alpha : \pi \circ s \to \pi \circ t$가 존재하여 도식
$$
\xymatrix{
\mathcal{S}_R \ar[r]_s \ar[d]_t & \mathcal{S}_U \ar[d]^\pi \\
\mathcal{S}_U \ar[r]^-\pi & [U/R]
}
$$
을 $2$-가환으로 만든다.
\end{lemma}

\begin{proof}
$S'$를 $S$ 위 스킴이라 하자. $r : S' \to R$를 $S$ 위 사상이라
하자. 그러면 $r \in R(S')$는 대상
$s \circ r, t \circ r \in U(S')$ 사이의 동형사상이다.
더욱이 이 구성은 당김과 양립한다. 이는 표준적인 $2$-사상
$\alpha_p : \pi_p \circ s \to \pi_p \circ t$를 준다.
여기서 $\pi_p : \mathcal{S}_U \to [U/_{\!p}R]$는 Lemma
\ref{spaces-groupoids-lemma-quotient-stack-arrows}의 증명에서와 같다. 따라서 도식
$$
\xymatrix{
\mathcal{S}_R \ar[r]_s \ar[d]^t & \mathcal{S}_U \ar[d]^{\pi_p} \\
\mathcal{S}_U \ar[r]^-{\pi_p} & [U/_{\!p}R]
}
$$
조차 $2$-가환이다. 그러므로 보조정리의 도식은 더욱이
$2$-가환이다.
\end{proof}

\begin{remark}
\label{spaces-groupoids-remark-fundamental-square}
뒤의 장들에서는 $\mathcal{S}_X$로 나타내는 집합에서의 스택이
$X$에 연관되어 있을 때 이를 단순히 $X$라고 쓰는 모호한
표기법을 사용할 것이다. 이 표기법을 사용하면 Lemma
\ref{spaces-groupoids-lemma-quotient-stack-2-arrow}의 도식은 익숙한 도식
$$
\xymatrix{
R \ar[r]_s \ar[d]_t & U \ar[d]^\pi \\
U \ar[r]^-\pi & [U/R]
}
$$
이 된다. 다음 절들에서는 이 도식이 좋은 성질을 많이 가짐을
보일 것이다. 특히 이것이 $2$-섬유곱임을
(Section \ref{spaces-groupoids-section-quotient-stack-2-cartesian}) 보이고,
$2$-쌍대등화자, 즉 $s$와 $t$를 쌍대등화하는 것에 가깝다는 것도
(Section \ref{spaces-groupoids-section-quotient-stacks-2-coequalize}) 보일 것이다.
\end{remark}




\section{몫 스택의 함자성}
\label{spaces-groupoids-section-functoriality-quotient-stacks}

\noindent
대수공간에서의 groupoid의 사상은 연관된 몫 스택의 사상을 준다.

\begin{lemma}
\label{spaces-groupoids-lemma-quotient-stack-functorial}
$S$를 스킴이라 하자. $B$를 $S$ 위 대수공간이라 하자.
$f : (U, R, s, t, c) \to (U', R', s', t', c')$를 $B$ 위
대수공간에서의 groupoid의 사상이라 하자. 그러면 $f$는 표준적인
몫 스택의 $1$-사상
$$
[f] : [U/R] \longrightarrow [U'/R'].
$$
을 유도한다.
\end{lemma}

\begin{proof}
함자 (\ref{spaces-groupoids-equation-quotient-stack})에 연관된 groupoid들로
섬유화된 범주를 $[U/_{\!p}R]$와 $[U'/_{\!p}R']$로 나타내자.
이들은 기저 자리 $(\Sch/S)_{fppf}$ 위에 있다.
$f$가 $1$-사상
$[U/_{\!p}R] \to [U'/_{\!p}R']$를 정의함은 명백하며, 이를
$[U'/R']$의 스택화 함자와 합성하여
$[U/_{\!p}R] \to [U'/R']$를 얻는다. 그런 다음 스택화 함자
$[U/_{\!p}R] \to [U/R]$의 보편성에 의해(Stacks, Lemma
\ref{stacks-lemma-stackify-groupoids-universal-property} 참조)
$[U/R] \to [U'/R']$를 얻는다.
\end{proof}

\noindent
Lemma \ref{spaces-groupoids-lemma-quotient-stack-functorial}에서와 같은
$B \to S$와
$f : (U, R, s, t, c) \to (U', R', s', t', c')$를 잡자.
이 상황에서 세 번째 $B$ 위 대수공간에서의 groupoid를 다음과 같이
정의한다. 여기서는 $T$-값 점의 언어를 사용하며, $T$는 달리는
스킴으로 $B$ 위에 있다.
\begin{enumerate}
\item $U'' = U \times_{f, U', t'} R'$로 둔다. 따라서 $T$-값 점은
쌍 $(u, r')$이며 $f(u) = t'(r')$를 만족한다.
\item $R'' = R \times_{f \circ s, U', t'} R'$로 둔다. 따라서
$T$-값 점은 쌍 $(r, r')$이며 $f(s(r)) = t'(r')$를 만족한다.
\item $s'' : R'' \to U''$는
$s''(r, r') = (s(r), r')$로 주어진다.
\item $t'' : R'' \to U''$는
$t''(r, r') = (t(r), c'(f(r), r'))$로 주어진다.
\item $c'' : R'' \times_{s'', U'', t''} R'' \to R''$는
$c''((r_1, r'_1), (r_2, r'_2)) = (c(r_1, r_2), r'_2)$로 주어진다.
\end{enumerate}
$c''$의 공식은 의미가 있다. 실제로
$s''(r_1, r'_1) = t''(r_2, r'_2)$이다.
$c''$가 결합적임은 명백하다. 항등원 $e''$는
$e''(u, r) = (e(u), r)$로 주어진다. $(r, r')$의 역원은
$(i(r), c'(f(r), r'))$로 주어진다. 따라서 실제로 $B$ 위
대수공간에서의 groupoid를 얻는다.

\medskip\noindent
사상 $U'' \to U$와 $R'' \to R$이 groupoid의 사상
$g : (U'', R'', s'', t'', c'') \to (U, R, s, t, c)$를
정의함은 명백하며, 이는 대수공간에서의 groupoid 사이에서 $B$
위에 있다. 더욱이 사상 $U'' \to U'$,
$(u, r') \mapsto s'(r')$와 $R'' \to U'$,
$(r, r') \mapsto s'(r')$는 실제로
$(U'', R'', s'', t'', c'')$가 $U'$ 위 대수공간에서의
groupoid임을 보인다.

\begin{lemma}
\label{spaces-groupoids-lemma-cartesian-square-of-morphism}
Lemma \ref{spaces-groupoids-lemma-quotient-stack-functorial}에서와 같은 표기법과
가정을 사용하자. 위에서 구성한 대수공간에서의 groupoid
$(U'', R'', s'', t'', c'')$를 $B$ 위에 잡자.
$2$-가환 정사각형
$$
\xymatrix{
[U''/R''] \ar[d] \ar[r]_{[g]} & [U/R] \ar[d]^{[f]} \\
\mathcal{S}_{U'} \ar[r] & [U'/R']
}
$$
이 존재하며, 이는 $[U''/R'']$를 $2$-섬유곱과 동일시한다.
\end{lemma}

\begin{proof}
사상 $[f]$와 $[g]$는 Lemma
\ref{spaces-groupoids-lemma-quotient-stack-functorial}을 적용하여 얻고, 나머지
두 사상은 Lemma \ref{spaces-groupoids-lemma-quotient-stack-arrows}에서 얻는다
($(U'', R'', s'', t'', c'')$가 $U'$ 위에 있다는 사실도 사용한다).
$2$-섬유곱 성질을 보이려면 groupoid들로 섬유화된 범주의 도식
$$
\xymatrix{
[U''/_{\!p}R''] \ar[d] \ar[r]_{[g]} & [U/_{\!p}R] \ar[d]^{[f]} \\
\mathcal{S}_{U'} \ar[r] & [U'/_{\!p}R']
}
$$
에 대해 보조정리를 증명하면 충분하다. Stacks, Lemma
\ref{stacks-lemma-stackification-fibre-product-categories-fibred-in-groupoids}를
보라. 달리 말하면, $2$-섬유곱
$\mathcal{S}_U \times_{[U'/_{\!p}R']} [U/_{\!p}R]$의 대상을
$T$ 위에서 잡았을 때, 이것이 $T$-값 점이며 $U''$에 속하는 점에
대응하고 사상에 대해서도 마찬가지임을 보이면 충분하다.
물론 이는 애초에 $U''$와 $R''$를 구성한 바로 그 방식이다.

\medskip\noindent
구체적으로,
$\mathcal{S}_U \times_{[U'/_{\!p}R']} [U/_{\!p}R]$의 대상은
$T$ 위에서 삼중항 $(u', u, r')$이다. 여기서 $u'$는 $T$-값 점이며
$U'$에 속하고, $u$는 $T$-값 점이며 $U$에 속한다. 또한
$r'$는 $u'$에서 $f(u)$로 가는 $[U'/R']_T$의 사상이다. 즉
$r'$는 $T$-값 점이며 $R$에 속하고
$s'(r') = u'$이고 $t'(r') = f(u)$이다. 정보를 잃지 않고
$u'$를 잊을 수 있음이 명백하고, 이 대상들이 $T$-값 점이며
$R''$에 속하는 점들과 일대일 대응함을 알 수 있다.

\medskip\noindent
사상도 마찬가지이다. $(u'_1, u_1, r'_1)$와
$(u'_2, u_2, r'_2)$를 $T$ 위 섬유곱의 두 대상이라 하자.
$(u'_2, u_2, r'_2)$에서 $(u'_1, u_1, r'_1)$로 가는 사상은
$(1, r)$로 주어진다. 여기서 $1 : u'_1 \to u'_2$는 단순히
$u'_1 = u'_2$를 뜻하고($\mathcal{S}_U$가 집합에서 섬유화되어
있기 때문이다), $r$은 $T$-값 점으로 $R$에 속하며
$s(r) = u_2$, $t(r) = u_1$ 및
$c'(f(r), r'_2) = r'_1$을 만족한다. 따라서 화살표
$$
(1, r) : (u'_2, u_2, r'_2) \to (u'_1, u_1, r'_1)
$$
는 쌍 $(r, r'_2)$만 알면 완전히 결정된다. 따라서 화살표의 함자는
$R''$로 표현되고, 더욱이 사상 $s''$, $t''$, $c''$는
$2$-섬유곱
$\mathcal{S}_U \times_{[U'/_{\!p}R']} [U/_{\!p}R]$에서의 시작점,
끝점, 합성에 각각 명백히 대응한다.
\end{proof}







\section{몫 스택의 2-데카르트 정사각형}
\label{spaces-groupoids-section-quotient-stack-2-cartesian}

\noindent
이 절에서는 몫 스택의 $\mathit{Isom}$-층들을 계산하고, 몫 스택을
정의하는 도식이 $2$-섬유곱임을 추론한다.

\begin{lemma}
\label{spaces-groupoids-lemma-quotient-stack-morphisms}
$B \to S$, $(U, R, s, t, c)$ 및
$\pi : \mathcal{S}_U \to [U/R]$가 Lemma
\ref{spaces-groupoids-lemma-quotient-stack-arrows}에서와 같다고 가정하자.
$S'$를 $S$ 위 스킴이라 하자. $x, y \in \Ob([U/R]_{S'})$를
$S'$ 위 몫 스택의 대상이라 하자. $x = \pi(x')$이고
$y = \pi(y')$라고 하자. 여기서 어떤 사상
$x', y' : S' \to U$가 존재한다고 가정한다. 그러면
$$
\mathit{Isom}(x, y) = S' \times_{(y', x'), U \times_S U} R
$$
이며, 이는 $S'$ 위 층의 등식이다.
\end{lemma}

\begin{proof}
$[U/_{\!p}R]$를 groupoid에서의 준층
(\ref{spaces-groupoids-equation-quotient-stack})에 연관된 groupoid들로 섬유화된
범주라 하자. 이는 Lemma \ref{spaces-groupoids-lemma-quotient-stack-arrows}의
증명에서와 같다. 구성상 층 $\mathit{Isom}(x, y)$는 준층
$\mathit{Isom}(x', y')$에 연관된 층이다. 한편
$[U/_{\!p}R]$에서의 사상의 정의에 의해
$$
\mathit{Isom}(x', y') = S' \times_{(y', x'), U \times_S U} R
$$
이고, 우변은 대수공간이므로 층이다.
\end{proof}

\begin{lemma}
\label{spaces-groupoids-lemma-quotient-stack-2-cartesian}
$B \to S$, $(U, R, s, t, c)$ 및
$\pi : \mathcal{S}_U \to [U/R]$가 Lemma
\ref{spaces-groupoids-lemma-quotient-stack-arrows}에서와 같다고 가정하자.
Lemma \ref{spaces-groupoids-lemma-quotient-stack-2-arrow}의 $2$-가환 정사각형
$$
\xymatrix{
\mathcal{S}_R \ar[r]_s \ar[d]_t & \mathcal{S}_U \ar[d]^\pi \\
\mathcal{S}_U \ar[r]^-\pi & [U/R]
}
$$
은 $2$-섬유곱이며, 이는 $(\Sch/S)_{fppf}$의 groupoid에서의
스택들 사이의 섬유곱이다.
\end{lemma}

\begin{proof}
Stacks, Lemma
\ref{stacks-lemma-2-product-stacks-in-groupoids}에 의하면 보조정리는
의미가 있다. 또한 그 결과에 의하면 함자
$$
\mathcal{S}_R \longrightarrow \mathcal{S}_U \times_{[U/R]} \mathcal{S}_U
$$
가 동치임을 보여야 한다. 이 함자는
$r : T \to R$를 $(T, t(r), s(r), \alpha(r))$로 보내고, 우변은
Categories, Lemma \ref{categories-lemma-2-product-categories-over-C}에
설명된 $2$-섬유곱이다. 정의들을 풀어 쓰면 이것은 정확히 Lemma
\ref{spaces-groupoids-lemma-quotient-stack-morphisms}의 내용이다. (대안적 증명:
이 상황에서 Lemma \ref{spaces-groupoids-lemma-cartesian-square-of-morphism}의 의미를
계산해도 결과를 얻는다.)
\end{proof}

\begin{lemma}
\label{spaces-groupoids-lemma-quotient-stack-isom}
$B \to S$와 $(U, R, s, t, c)$가 Definition
\ref{spaces-groupoids-definition-quotient-stack} (1)에서와 같다고 가정하자.
$T$를 $S$ 위 임의의 스킴이라 하고, 대상 $x, y$를 $[U/R]$에서
$T$ 위에 잡자. 층 $\mathit{Isom}(x, y)$는
$(\Sch/T)_{fppf}$ 위에서 다음
성질을 갖는다. fppf 덮개 $\{T_i \to T\}_{i \in I}$가 존재하여
$$
\mathit{Isom}(x, y)|_{(\Sch/T_i)_{fppf}}
$$
는 대수공간으로 표현된다.
\end{lemma}

\begin{proof}
Lemma \ref{spaces-groupoids-lemma-quotient-stack-morphisms}와 다음 사실에서 바로
따른다. 몫 스택의 구성상 $x$와 $y$는 모두 fppf 위상에서
국소적으로 $\mathcal{S}_U$의 대상에서 온다.
\end{proof}


























\section{몫 스택의 2-쌍대등화자 성질}
\label{spaces-groupoids-section-quotient-stacks-2-coequalize}

\noindent
groupoid에는 합성이 있으며, 이는 위 보조정리의 표준적인
$2$-사상에 대한 여접합 조건을 낳는다. 정확히 정식화하기 위해
Categories, Sections \ref{categories-section-formal-cat-cat}와
\ref{categories-section-2-categories}에서 도입한 표기법을
사용할 것이다.

\begin{lemma}
\label{spaces-groupoids-lemma-quotient-stack-cocycle}
Lemmas \ref{spaces-groupoids-lemma-quotient-stack-arrows}와
\ref{spaces-groupoids-lemma-quotient-stack-2-arrow}에서와 같은 가정과 표기법을
사용하자. 다음의 수직 합성
$$
\xymatrix@C=15pc{
\mathcal{S}_{R \times_{s, U, t} R}
\ruppertwocell^{\pi \circ s \circ \text{pr}_1 = \pi \circ s \circ c}{\ \ \ \ \ \ \alpha \star \text{id}_{\text{pr}_1}}
\ar[r]_(.3){\pi \circ t \circ \text{pr}_1 = \pi \circ s \circ \text{pr}_0}
\rlowertwocell_{\pi \circ t \circ \text{pr}_0 = \pi \circ t \circ c}{\ \ \ \ \ \ \alpha \star \text{id}_{\text{pr}_0}}
&
[U/R]
}
$$
은 $2$-사상 $\alpha \star \text{id}_c$이다. 공식으로는
$\alpha \star \text{id}_c =
(\alpha \star \text{id}_{\text{pr}_0})
\circ
(\alpha \star \text{id}_{\text{pr}_1})
$이다.
\end{lemma}

\begin{proof}
두 가지를 지적한다.
\begin{enumerate}
\item 공식
$\alpha \star \text{id}_c = (\alpha \star \text{id}_{\text{pr}_0}) \circ
(\alpha \star \text{id}_{\text{pr}_1})$이 의미를 가지려면 다음
{\it 등식들}을 알아야 한다.
$\pi \circ s \circ \text{pr}_1 = \pi \circ s \circ c$,
$\pi \circ t \circ \text{pr}_1 = \pi \circ s \circ \text{pr}_0$,
그리고
$\pi \circ t \circ \text{pr}_0 = \pi \circ t \circ c$이다.
즉 두 번째 등식은 수직 합성 $\circ$가 의미를 갖게 하고, 나머지
두 등식은 공식의 양변이 시작점과 끝점이 같은 $2$-사상임을
보장한다.
\item 이 보조정리가 성립하는 이유는 groupoid에서의 준층
(\ref{spaces-groupoids-equation-quotient-stack})에 연관된 groupoid들로 섬유화된
범주 $[U/_{\!p}R]$에서의 합성이 합성법칙
$c : R \times_{s, U, t} R \to R$에서 오기 때문이다.
\end{enumerate}
보조정리의 증명은 생략한다.
\end{proof}

\noindent
보조정리의 상황에서는 실제로 등식
$s \circ \text{pr}_1 = s \circ c$,
$t \circ \text{pr}_1 = s \circ \text{pr}_0$, 그리고
$t \circ \text{pr}_0 = t \circ c$가 있으며, 이는 $\pi$와
합성하기 전부터 성립함에 유의하라.
따라서 아래 보조정리의 공식은 위 보조정리의 공식과 정확히 같은
방식으로 의미를 갖는다.

\begin{lemma}
\label{spaces-groupoids-lemma-quotient-stack-2-coequalizer}
Lemmas \ref{spaces-groupoids-lemma-quotient-stack-arrows}와
\ref{spaces-groupoids-lemma-quotient-stack-2-arrow}에서와 같은 가정과 표기법을
사용하자. Lemma \ref{spaces-groupoids-lemma-quotient-stack-2-arrow}의
$2$-가환 도식은 다음 의미에서 $2$-쌍대등화자이다. 다음 자료가
주어졌다고 하자.
\begin{enumerate}
\item groupoid에서의 스택 $\mathcal{X}$, 이는
$(\Sch/S)_{fppf}$ 위에 있다.
\item $1$-사상 $f : \mathcal{S}_U \to \mathcal{X}$,
\item $2$-화살표 $\beta : f \circ s \to f \circ t$.
\end{enumerate}
이들이
$$
\beta \star \text{id}_c
=
(\beta \star \text{id}_{\text{pr}_0})
\circ
(\beta \star \text{id}_{\text{pr}_1})
$$
를 만족하면, $1$-사상 $[U/R] \to \mathcal{X}$가 존재하여 도식
$$
\xymatrix{
\mathcal{S}_R \ar[r]_s \ar[d]^t & \mathcal{S}_U \ar[d] \ar[ddr]^f \\
\mathcal{S}_U \ar[r] \ar[rrd]_f & [U/R] \ar[rd] \\
& & \mathcal{X}
}
$$
을 $2$-가환으로 만든다.
\end{lemma}

\begin{proof}
보조정리에서와 같은 $\mathcal{X}$, $f$, $\beta$가 주어졌다고
하자. Stacks, Lemma
\ref{stacks-lemma-stackify-groupoids-universal-property}에 의해
$1$-사상 $g : [U/_{\!p}R] \to \mathcal{X}$를 구성하면 충분하다.
먼저 $1$-사상 $\mathcal{S}_U \to [U/_{\!p}R]$가 대상들 위에서
전단사임에 유의하라. 따라서 대상들 위에서는
$g(x) = f(x)$로 둘 수 있다. 여기서
$x \in \Ob(\mathcal{S}_U) = \Ob([U/_{\!p}R])$이다.
사상 $\varphi : x \to y$는 $[U/_{\!p}R]$에 속하며 가환 도식
$$
\xymatrix{
S_2 \ar[dd]_h \ar[r]_x \ar[dr]_\varphi & U \\
& R \ar[u]_s \ar[d]^t \\
S_1 \ar[r]^y & U.
}
$$
에서 온다. 따라서 $g(\varphi)$를 합성
$$
\xymatrix{
f(x) \ar@{=}[r] \ar[rrrrrd] &
f(s \circ \varphi) \ar@{=}[r] &
(f \circ s)(\varphi) \ar[r]^\beta &
(f \circ t)(\varphi) \ar@{=}[r] &
f(t \circ \varphi) \ar@{=}[r] &
f(y \circ h) \ar[d] \\
& & & & & f(y).
}
$$
과 같게 둘 수 있다. 수직 화살표는 함자 $f$를 표준적인 사상
$y \circ h \to y$에 적용한 결과이다. 이 사상은
$\mathcal{S}_U$에 속하며, $h$를 올리고 끝점이 $y$인 강한
데카르트 사상이다.
이렇게 정의한 $f$가 적어도 섬유 범주들 위에서 합성과
양립함을 확인하자. $S'$를 $S$ 위 스킴이라 하고
$a : S' \to R \times_{s, U, t} R$를 사상이라 하자.
이 상황에서
$x = s \circ \text{pr}_1 \circ a = s \circ c \circ a$,
$y = t \circ \text{pr}_1 \circ a = s \circ \text{pr}_0 \circ a$,
그리고
$z = t \circ \text{pr}_0 \circ a = t \circ \text{pr}_0 \circ c$로
두어 가환 도식
$$
\xymatrix{
x \ar[rr]_{c \circ a} \ar[rd]_{\text{pr}_1 \circ a} & & z \\
& y \ar[ru]_{\text{pr}_0 \circ a}
}
$$
을 섬유 범주 $[U/_{\!p}R]_{S'}$에서 얻는다.
더욱이 이 섬유 범주의 모든 가환 삼각형은 이러한 꼴이다.
위 정의들에 의해 $f$가 이를 가환 도식으로 보낼 필요충분조건은
도식
$$
\xymatrix{
& (f \circ s)(c \circ a) \ar[r]_-{\beta} &
(f \circ t)(c \circ a) \ar@{=}[rd] & \\
(f \circ s)(\text{pr}_1 \circ a) \ar[rd]^\beta \ar@{=}[ru] & & &
(f \circ t)(\text{pr}_0 \circ a) \\
& (f \circ t)(\text{pr}_1 \circ a) \ar@{=}[r] &
(f \circ s)(\text{pr}_0 \circ a) \ar[ru]^\beta
}
$$
이 가환하는 것이며, 이는 정확히 보조정리의 공식이 나타내는
조건이다. $f$가 항등원을 항등원으로 보내고 임의의 사상들의
합성과 양립한다는 확인은 생략한다.
\end{proof}













\section{몫 스택의 명시적 설명}
\label{spaces-groupoids-section-explicit-quotient-stacks}

\noindent
결과를 정식화하려면 몇 가지 표기법을 도입해야 한다.
$B \to S$와 $(U, R, s, t, c)$가 Definition
\ref{spaces-groupoids-definition-quotient-stack} (1)에서와 같다고 가정하자.
$T$를 $S$ 위 스킴이라 하자.
$\mathcal{T} = \{T_i \to T\}_{i \in I}$를 fppf 덮개라 하자.
{\it $[U/R]$-하강 자료} 가운데 $\mathcal{T}$에 대한 것은 계
$(u_i, r_{ij})$로 주어지는데, 여기서
\begin{enumerate}
\item 각 $i$에 대하여 사상 $u_i : T_i \to U$가 있고,
\item 각 $i, j$에 대하여 사상
$r_{ij} : T_i \times_T T_j \to R$가 있으며,
\end{enumerate}
다음을 만족한다.
\begin{enumerate}
\item[(a)] $T_i \times_T T_j \to U$의 사상으로서
$$
s \circ r_{ij} = u_i \circ \text{pr}_0
\quad\text{and}\quad
t \circ r_{ij} = u_j \circ \text{pr}_1,
$$
이고,
\item[(b)] $T_i \times_T T_j \times_T T_k \to R$의 사상으로서
$$
c \circ (r_{jk} \circ \text{pr}_{12}, r_{ij} \circ \text{pr}_{01})
= r_{ik} \circ \text{pr}_{02}.
$$
이다.
\end{enumerate}
두 하강 자료 사이의 {\it 사상
$(u_i, r_{ij}) \to (u'_i, r'_{ij})$} 가운데 $[U/R]$-하강
자료의 사상으로서 같은 덮개 $\mathcal{T}$ 위에 있는 것은 다음을 만족하는
모음 $(r_i : T_i \to R)$이다.
\begin{enumerate}
\item[] $(\alpha)$ $T_i \to U$의 사상으로서
$$
u_i = s \circ r_i
\quad\text{and}\quad
u'_i = t \circ r_i
$$
이고,
\item[] $(\beta)$ $T_i \times_T T_j \to R$의 사상으로서
$$
c \circ (r'_{ij}, r_i \circ \text{pr}_0)
=
c \circ (r_j \circ \text{pr}_1, r_{ij}).
$$
이다.
\end{enumerate}
고정된 덮개에 대한 하강 자료의 사상들 위에는 자연스러운 합성법칙이
있고, 하강 자료의 범주를 얻는다. 이 범주는 groupoid이다.
마지막으로
$\mathcal{T}' = \{T'_j \to T\}_{j \in J}$가
$\mathcal{T}$를 세분하는 두 번째 fppf 덮개이면 하강 자료의 당김
개념이 있다. 이 경우 이를 명시적으로 설명하기가 특히 쉽다.
즉 $\alpha : J \to I$이고
% P09-ERR-1458
$\varphi_j : T'_j \to T_{\alpha(i)}$가 덮개의 사상이면 하강 자료
$(u_i, r_{ii'})$의 당김은 단순히
$$
(u_{\alpha(i)} \circ \varphi_j,
r_{\alpha(j)\alpha(j')} \circ \varphi_j \times \varphi_{j'}).
$$
이다. 이렇게 정의한 당김은 $\mathcal{T}$ 위 하강 자료의 범주에서
$\mathcal{T}'$ 위 하강 자료의 범주로 가는 함자를 정의한다.

\begin{lemma}
\label{spaces-groupoids-lemma-quotient-stack-objects}
$B \to S$와 $(U, R, s, t, c)$가 Definition
\ref{spaces-groupoids-definition-quotient-stack} (1)에서와 같다고 가정하자.
$\pi : \mathcal{S}_U \to [U/R]$를 Lemma
\ref{spaces-groupoids-lemma-quotient-stack-arrows}에서와 같이 잡자.
$T$를 $S$ 위 스킴이라 하자.
\begin{enumerate}
\item 모든 대상 $x$가 섬유 범주 $[U/R]_T$에 속하면, fppf 덮개
$\{f_i : T_i \to T\}_{i \in I}$가 존재하여
$f_i^*x \cong \pi(u_i)$가 되며, 여기서 어떤
$u_i \in U(T_i)$를 취할 수 있다.
\item 동형사상들의 합성
$$
\pi(u_i \circ \text{pr}_0)
=
\text{pr}_0^*\pi(u_i)
\cong
\text{pr}_0^*f_i^*x
\cong
\text{pr}_1^*f_j^*x
\cong
\text{pr}_1^*\pi(u_j)
=
\pi(u_j \circ \text{pr}_1)
$$
은 $\pi(r_{ij})$ 꼴이며, 여기서 어떤 사상
$r_{ij} : T_i \times_T T_j \to R$를 취할 수 있다.
\item 계 $(u_i, r_{ij})$는 위에서 정의한
$[U/R]$-하강 자료를 이룬다.
\item 임의의 $[U/R]$-하강 자료 $(u_i, r_{ij})$는 이 방식으로
생긴다.
\item $x$가 위와 같이 $(u_i, r_{ij})$에 대응하고,
$y \in \Ob([U/R]_T)$가 $(u'_i, r'_{ij})$에 대응하면 표준적인
전단사
$$
\Mor_{[U/R]_T}(x, y)
\longleftrightarrow
\left\{
\begin{matrix}
\text{morphisms }(u_i, r_{ij}) \to (u'_i, r'_{ij})\\
\text{of }[U/R]\text{-descent data}
\end{matrix}
\right\}
$$
가 존재한다.
\item 이 대응은 fppf 덮개의 세분과 양립한다.
\end{enumerate}
\end{lemma}

\begin{proof}
명제 (1)은 스택화 구성의 일부이다. (2)는 Lemma
\ref{spaces-groupoids-lemma-quotient-stack-morphisms}에서 따른다.
(3)의 확인은 생략한다. (4)는 스택에서 모든 하강 자료가
효과적이라는 사실을 옮겨 쓴 것이다. (5)와 (6)의 확인은 생략한다.
\end{proof}







\section{제한과 몫 스택}
\label{spaces-groupoids-section-quotient-stack-restrict}

\noindent
이 절에서는 제한을 취할 때 몫 스택에 어떤 일이 일어나는지 연구한다.

\begin{lemma}
\label{spaces-groupoids-lemma-quotient-stack-restrict}
Lemma \ref{spaces-groupoids-lemma-quotient-stack-functorial}에서와 같은 표기법과
가정을 사용하자. 몫 스택의 사상
$$
[f] : [U/R] \longrightarrow [U'/R']
$$
이 완전충실일 필요충분조건은 $R$가 $R'$의 제한인 것이다.
여기서 제한은 사상 $f : U \to U'$를 통하여 취한다.
\end{lemma}

\begin{proof}
$x, y$를 $[U/R]$의 대상으로 하되 스킴 $T/S$ 위에 잡자.
$x', y'$를 $x, y$의 상이라 하되 범주는 $[U'/R']_T$라 하자.
함자 $[f]$가 완전충실일 필요충분조건은 층의 사상
$$
\mathit{Isom}(x, y) \longrightarrow \mathit{Isom}(x', y')
$$
이 모든 $T, x, y$에 대하여 동형사상인 것이다. 이를 $T$ 위에서
국소적으로(fppf 위상으로)
검사할 수 있다. 따라서 Lemma \ref{spaces-groupoids-lemma-quotient-stack-objects}에
의해 $x, y$가 $a, b \in U(T)$에서 온다고 가정할 수 있다.
그 경우 $x', y'$가 $f \circ a, f \circ b$에 대응함을 안다.
Lemma \ref{spaces-groupoids-lemma-quotient-stack-morphisms}에 의해 이 경우 표시된
층의 사상은
$$
T \times_{(a, b), U \times_B U} R
\longrightarrow
T \times_{f \circ a, f \circ b, U' \times_B U'} R'.
$$
이 된다. $R$가 제한이면 이것은 동형사상이다. 실제로 그 경우
$R = (U \times_B U) \times_{U' \times_B U'} R'$이다.
Lemma \ref{spaces-groupoids-lemma-restrict-groupoid-relation}과 그 증명을 보라.
거꾸로 마지막 표시 사상이 모든 $T, a, b$에 대하여 동형사상이면
$R = (U \times_B U) \times_{U' \times_B U'} R'$가 따르며, 즉
$R$는 $R'$의 제한이다.
\end{proof}

\begin{lemma}
\label{spaces-groupoids-lemma-quotient-stack-restrict-equivalence}
Lemma \ref{spaces-groupoids-lemma-quotient-stack-functorial}에서와 같은 표기법과
가정을 사용하자. 몫 스택의 사상
$$
[f] : [U/R] \longrightarrow [U'/R']
$$
이 동치일 필요충분조건은 다음과 같다.
\begin{enumerate}
\item $(U, R, s, t, c)$는 $(U', R', s', t', c')$의 제한이며,
이는 $f : U \to U'$를 통하여 취한다.
\item 사상
$$
\xymatrix{
U \times_{f, U', t'} R' \ar[r]_-{\text{pr}_1} \ar@/^3ex/[rr]^h
& R' \ar[r]_{s'} & U'
}
$$
은 층의 전사이다.
\end{enumerate}
예를 들어 $\{h : U \times_{f, U', t'} R' \to U'\}$가 fppf
덮개이거나, $f : U \to U'$가 층의 전사이거나,
$\{f : U \to U'\}$가 fppf 덮개이면 (2)가 성립한다.
\end{lemma}

\begin{proof}
Lemma \ref{spaces-groupoids-lemma-quotient-stack-restrict}에 의해 (1)이 완전충실성과
동치임을 이미 안다. 따라서 (1)이 성립하고 $[f]$가 완전충실이라고
가정할 수 있다. 이 가정 아래 $[f]$가 동치일 필요충분조건이 (2)임을
보이는 것이 목표이다. 동치를 특징짓는 Stacks, Lemma
\ref{stacks-lemma-characterize-essentially-surjective-when-ff}를
사용할 수 있다.

\medskip\noindent
(2)를 가정하자. Stacks, Lemma
\ref{stacks-lemma-characterize-essentially-surjective-when-ff}를
사용하여 $[f]$가 동치임을 증명하겠다.
$T$가 스킴이고 $x' \in \Ob([U'/R']_T)$라고 하자.
$\{g_i : T_i \to T\}$라는 덮개가 존재하여 $g_i^*x'$가 어떤
원소 $a'_i \in U'(T_i)$의 상이 된다. Lemma
\ref{spaces-groupoids-lemma-quotient-stack-objects}를 보라. 따라서 $x'$가
$a' \in U'(T)$의 상이라고 가정할 수 있다.
$h$가 층의 전사라는 가정에 의해 fppf 덮개
$\{\varphi_i : T_i \to T\}$와 다음을 만족하는 사상
% P09-ERR-1459
$b_i : T_i \to U \times_{g, U', t'} R'$를 찾을 수 있다.
$a' \circ \varphi_i = h \circ b_i$이다.
$a_i = \text{pr}_0 \circ b_i : T_i \to U$라 나타내자.
그러면 $a_i \in U(T_i)$는 $f \circ a_i \in U'(T_i)$로 사상되고
$f \circ a_i \cong_{T_i} h \circ b_i = a' \circ \varphi_i$이다.
여기서 $\cong_{T_i}$는 섬유 범주 $[U'/R']_{T_i}$에서의
동형사상을 나타낸다. 즉 동형사상을 주는 $R'(T_i)$의 원소는
$\text{pr}_1 \circ b_i$이다. 이는 원하는 대로 $x$의 $T_i$로의
제한이 함자 $[U/R]_{T_i} \to [U'/R']_{T_i}$의 본질적 상에 있음을
뜻한다.

\medskip\noindent
$[f]$가 동치라고 가정하자.
$\xi' \in [U'/R']_{U'}$를 $U'$의 항등사상에 대응하는 대상이라
하자. Stacks, Lemma
\ref{stacks-lemma-characterize-essentially-surjective-when-ff}를
적용하면 fppf 덮개
$\mathcal{U}' = \{g'_i : U'_i \to U'\}$가 존재하여
$(g'_i)^*\xi' \cong [f](\xi_i)$가 된다. 여기서 어떤
$\xi_i$가 $[U/R]_{U'_i}$에 있다.
덮개 $\mathcal{U}'$를 세분한 뒤(Lemma
\ref{spaces-groupoids-lemma-quotient-stack-objects} 사용) $\xi_i$가 사상
$a_i : U'_i \to U$에서 온다고 가정할 수 있다.
$[f](\xi_i) \cong (g'_i)^*\xi'$라는 사실은, 필요하면 덮개
$\mathcal{U}'$를 한 번 더 세분한 뒤, 다음을 만족하는 사상
$r'_i : U'_i \to R'$가 존재한다는 뜻이다.
$t' \circ r'_i = f \circ a_i$이고
$s' \circ r'_i = \text{id}_{U'} \circ g'_i$이다. 도식은
$$
\xymatrix{
U \ar[d]^f & & U'_i \ar[ll]^{a_i} \ar[ld]_{r'_i} \ar[d]^{g'_i} \\
U' & R' \ar[l]_{t'} \ar[r]^{s'} & U'
}
$$
이다. 따라서
% P09-ERR-1460
$(a_i, r'_i) : U'_i \to U \times_{g, U', t'} R'$는
$h \circ (a_i, r'_i) = g'_i$를 만족하는 사상이다. 그러므로
$\{h : U \times_{g, U', t'} R' \to U'\}$는 fppf 덮개
$\mathcal{U}'$로 세분될 수 있고, 이는 $h$가 층의 전사를
유도한다는 뜻이다. Topologies on Spaces, Lemma
\ref{spaces-topologies-lemma-fppf-covering-surjective}를 보라.

\medskip\noindent
$\{h\}$가 fppf 덮개이면 층의 전사를 유도한다. Topologies on
Spaces, Lemma
\ref{spaces-topologies-lemma-fppf-covering-surjective}를 보라.
% P09-ERR-1461
$U' \to U$가 전사이면 $h$도 전사이다. 실제로 $s$가 단면,
즉 대수공간에서의 groupoid의 중립원 $e$를 가진다.
\end{proof}

\begin{lemma}
\label{spaces-groupoids-lemma-criterion-fibre-product}
Lemma \ref{spaces-groupoids-lemma-quotient-stack-functorial}에서와 같은 표기법과
가정을 사용하자. 다음 도식이 데카르트라고 가정하자.
$$
\xymatrix{
R \ar[d]_s \ar[r]_f & R' \ar[d]^{s'} \\
U \ar[r]^f & U'
}
$$
그러면
$$
\xymatrix{
\mathcal{S}_U \ar[d] \ar[r] & [U/R] \ar[d]^{[f]} \\
\mathcal{S}_{U'} \ar[r] & [U'/R']
}
$$
은 $2$-섬유곱 정사각형이다.
\end{lemma}

\begin{proof}
보조정리의 명제에 나온 (첫 번째) 데카르트 도식에 역동형사상
$i : R \to R$와 $i' : R' \to R'$를 적용하면
$$
\xymatrix{
R \ar[d]_t \ar[r]_f & R' \ar[d]^{t'} \\
U \ar[r]^f & U'
}
$$
도 데카르트임을 안다. Lemma
\ref{spaces-groupoids-lemma-cartesian-square-of-morphism}에 의해 $2$-섬유
정사각형
$$
\xymatrix{
[U''/R''] \ar[d] \ar[r] & [U/R] \ar[d] \\
\mathcal{S}_{U'} \ar[r] & [U'/R']
}
$$
을 갖는다. 여기서
$U'' = U \times_{f, U', t'} R'$이고
$R'' = R \times_{f \circ s, U', t'} R'$이다.
위에서 $(t, f) : R \to U''$가 동형사상이고
$$
R'' =
R \times_{f \circ s, U', t'} R' =
R \times_{s, U} U \times_{f, U', t'} R' =
% P09-ERR-1462
R \times_{s, U, t} \times R.
$$
임을 안다. 명시적으로 동형사상
$R \times_{s, U, t} R \to R''$은 규칙
$(r_0, r_1) \mapsto (r_0, f(r_1))$로 주어진다.
더욱이 $s'', t'', c''$는 사상
$$
R \times_{s, U, t} R \to R,
\quad
s''(r_0, r_1) = r_1, \quad t''(r_0, r_1) = c(r_0, r_1)
$$
및
$$
\begin{matrix}
c'' : &
(R \times_{s, U, t} R) \times_{s'', R, t''} (R \times_{s, U, t} R)
&
\longrightarrow
&
R \times_{s, U, t} R, \\
&
((r_0, r_1), (r_2, r_3)) &
\longmapsto & (c(r_0, r_2), r_3).
\end{matrix}
$$
으로 옮겨진다. 동형사상
$$
R \times_{s, U, s} R \longrightarrow R \times_{s, U, t} R,
\quad
(r_0, r_1) \longmapsto (c(r_0, i(r_1)), r_1)
$$
을 앞에 합성하면 $t''$와 $s''$는
$\text{pr}_0$와 $\text{pr}_1$로, $c''$는
$\text{pr}_{02} :
R \times_{s, U, s} R \times_{s, U, s} R \to R \times_{s, U, s} R$로
바뀐다. 따라서 동형사상
$[U''/R''] \cong [R/R \times_{s, U, s} R]$가 존재함을 안다.
여기서 대수공간에서의 groupoid
$(R, R \times_{s, U, s} R, s'', t'', c'')$는 자명한 groupoid
$(U, U, \text{id}, \text{id}, \text{id})$를
$s : R \to U$를 통하여 제한한 것이다.
$s : R \to U$는 오른쪽 역을 가지므로 fppf 층의 전사이고, 따라서
사상
$$
[U''/R''] \cong [R/R \times_{s, U, s} R]
\longrightarrow
[U/U] = \mathcal{S}_U
$$
은 Lemma \ref{spaces-groupoids-lemma-quotient-stack-restrict-equivalence}에 의해
동치이다. 이것으로 보조정리를 증명했다.
\end{proof}





\section{관성과 몫 스택}
\label{spaces-groupoids-section-inertia}

\noindent
groupoid에서의 스택의 (상대적) 관성 스택은 Stacks, Section
\ref{stacks-section-the-inertia-stack}에서 정의한다.
섬유화된 범주의 설정에서 실제 구성과 그 성질 몇 가지는 Categories,
Section \ref{categories-section-inertia}에 있다.

\begin{lemma}
\label{spaces-groupoids-lemma-presentation-inertia}
$B \to S$와 $(U, R, s, t, c)$가 Definition
\ref{spaces-groupoids-definition-quotient-stack} (1)에서와 같다고 가정하자.
$G/U$를 groupoid $(U, R, s, t, c, e, i)$의 안정자 군
대수공간이라 하자. Definition
\ref{spaces-groupoids-definition-stabilizer-groupoid}을 보라.
$R' = R \times_{s, U} G$로 두고 다음과 같이 둔다.
\begin{enumerate}
\item $s' : R' \to G$, $(r, g) \mapsto g$,
\item $t' : R' \to G$, $(r, g) \mapsto c(r, c(g, i(r)))$,
\item $c' : R' \times_{s', G, t'} R' \to R'$,
$((r_1, g_1), (r_2, g_2) \mapsto (c(r_1, r_2), g_1)$.
\end{enumerate}
그러면 $(G, R', s', t', c')$는 $B$ 위 대수공간에서의
groupoid이고
$$
\mathcal{I}_{[U/R]} = [G/ R'].
$$
이다. 즉 연관된 몫 스택은 $[U/R]$의 관성 스택이다.
\end{lemma}

\begin{proof}
Stacks, Lemma \ref{stacks-lemma-stackification-inertia}에 의해
$\mathcal{I}_{[U/_{\!p}R]} = [G/_{\!p} R']$임을 증명하면 충분하다.
$T$를 $S$ 위 스킴이라 하자. $[U/_{\!p}R]$의 관성
섬유화 범주에서 $T$ 위의 대상은 쌍 $(x, g)$로 주어짐을 상기하라.
여기서 $x$는 $[U/_{\!\!p}R]$의 대상으로서 $T$ 위에 있고,
$g$는 $x$의 자기동형사상으로서 그 섬유 범주에서 $T$ 위에 있다.
달리 말하면 $x : T \to U$와 $g : T \to R$이고
$x = s \circ g = t \circ g$를 만족한다. 이는 정확히
$g : T \to G$라는 뜻이다. 관성 섬유화 범주에서
$(x, g) \to (y, h)$로 가는 $T$ 위 사상은
$r : T \to R$로 주어지고, $s(r) = x$, $t(r) = y$,
$c(r, g) = c(h, r)$를 만족한다. Categories, Lemma
\ref{categories-lemma-inertia-fibred-category}의 가환 도식을
보라. 공식으로는
$$
h = c(r, c(g, i(r))) = c(c(r, g), i(r)).
$$
이다. 표기법 $s(r)$ 등은 $s \circ r$ 등의 약기이다.
$r_1 : (x_2, g_2) \to (x_1, g_1)$과
$r_2 : (x_1, g_1) \to (x_2, g_2)$의 합성은
% P09-ERR-1463
$c(r_1, r_2) : (x_1, g_1) \to (x_3, g_3)$이다.

\medskip\noindent
위에서 대상을 나타낼 때 $g$라고 쓸 수도 있었다. 이는 $(x, g)$를
대신하는 표기이다.
이 대상은 $\mathcal{I}_{[U/_{\!p}R]}$에서 $T$ 위에 있다.
실제로 $x$는 $g$의 상이며, 그 사상은 $G \to U$이다.
그러면 사상 $g \to h$는 $\mathcal{I}_{[U/_{\!p}R]}$에서
$T$ 위에 있고, 사상 $r' : T \to R'$와 정확히 대응한다.
이때 $s'(r') = g$, $t'(r') = h$이다. 더욱이 합성은 (3)에서
설명한 규칙에 대응한다. 따라서 보조정리가 증명되었다.
\end{proof}

\begin{lemma}
\label{spaces-groupoids-lemma-2-cartesian-inertia}
$B \to S$와 $(U, R, s, t, c)$가 Definition
\ref{spaces-groupoids-definition-quotient-stack} (1)에서와 같다고 가정하자.
$G/U$를 groupoid $(U, R, s, t, c, e, i)$의 안정자 군
대수공간이라 하자. Definition
\ref{spaces-groupoids-definition-stabilizer-groupoid}을 보라.
표준적인 $2$-데카르트 도식
$$
\xymatrix{
\mathcal{S}_G \ar[r] \ar[d] & \mathcal{S}_U \ar[d] \\
\mathcal{I}_{[U/R]} \ar[r] & [U/R]
}
$$
이 존재하며, 이는 $(\Sch/S)_{fppf}$의 groupoid에서의
스택들의 도식이다.
\end{lemma}

\begin{proof}
Lemma \ref{spaces-groupoids-lemma-criterion-fibre-product}에 의해 Lemma
\ref{spaces-groupoids-lemma-presentation-inertia}의 사상 $s' : R' \to G$가
$s$의 기저변환과 동형임을 증명하면 충분하다. 이 기저변환은 구조
사상 $G \to U$에 의해 취한다.
이 기저변환 성질은 $s'$의 구성에서 명백하다.
\end{proof}






\section{제르브와 몫 스택}
\label{spaces-groupoids-section-gerbes}

\noindent
이 절에서는 몫 스택을 Stacks, Section
\ref{stacks-section-gerbes}의 논의, 특히 Stacks, Definition
\ref{stacks-definition-gerbe-over-stack-in-groupoids}에서 정의한
제르브와 관련짓는다. 이 절에 등장하는 groupoid에서의 스택들은
일반적으로 대수 스택이 아니다!

\begin{lemma}
\label{spaces-groupoids-lemma-when-gerbe}
Lemma \ref{spaces-groupoids-lemma-quotient-stack-functorial}에서와 같은 표기법과
가정을 사용하자. 몫 스택의 사상
$$
[f] : [U/R] \longrightarrow [U'/R']
$$
은 $[U/R]$를 $[U'/R']$ 위 제르브로 만든다. 단,
$f : U \to U'$와 $R \to R'|_U$가 fppf 층의 전사라고
가정한다. 여기서 $R'|_U$는 $R'$의 $U$로의 제한이며,
$f : U \to U'$를 통하여 취한다.
\end{lemma}

\begin{proof}
Stacks, Lemma \ref{stacks-lemma-when-gerbe}의 성질 (2) (a)와
(2) (b)가 성립함을 확인하겠다. $U \to U'$가 층의 전사이므로
성질 (2)(a)가 성립한다($[U'/R']$의 대상들이 국소적으로 $U'$에서
옴을 보려면 Lemma \ref{spaces-groupoids-lemma-quotient-stack-objects}를 사용하라).
(2)(b)를 증명하기 위해 $x, y$를 $[U/R]$의 대상으로 하되
스킴 $T/S$ 위에 잡자. $x', y'$를 $x, y$의 상이라 하되 범주는
% P09-ERR-1464
$[U'/'R]_T$라 하자. 조건 (2)(b)는 층의 사상
$$
\mathit{Isom}(x, y) \longrightarrow \mathit{Isom}(x', y')
$$
이 $(\Sch/T)_{fppf}$ 위에서 전사임을 확인하라고 요구한다.
이를 보이기 위해 $T$ 위에서
fppf 국소적으로 작업할 수 있고,
% P09-ERR-1465
그 대상들이 $a, b \in U(T)$에서 온다고 가정할 수 있다.
그 경우 $x', y'$는 $f \circ a, f \circ b$에 대응한다.
Lemma \ref{spaces-groupoids-lemma-quotient-stack-morphisms}에 의해 이 경우 표시한
층의 사상은
$$
T \times_{(a, b), U \times_B U} R
\longrightarrow
T \times_{f \circ a, f \circ b, U' \times_B U'} R' =
T \times_{(a, b), U \times_B U} R'|_U.
$$
이 된다. 따라서 $R \to R'|_U$가 fppf 층의 전사라는 가정은
$(\Sch/S)_{fppf}$ 위에서 원하는 전사성을
함의한다.
\end{proof}

\begin{lemma}
\label{spaces-groupoids-lemma-group-quotient-gerbe}
$S$를 스킴이라 하자. $B$를 $S$ 위 대수공간이라 하자.
$G$를 $B$ 위 군 대수공간이라 하자. $B$에 $G$의 자명한 작용을
주자. 사상
$$
[B/G] \longrightarrow \mathcal{S}_B
$$
(Lemma \ref{spaces-groupoids-lemma-quotient-stack-arrows})는 $[B/G]$를 $B$ 위
제르브로 만든다.
\end{lemma}

\begin{proof}
사상 $B \to B$와 $B \times_B G \to B$가 층의 사상으로서
전사이므로 Lemma \ref{spaces-groupoids-lemma-when-gerbe}에서 바로 따른다.
\end{proof}








\section{몫 스택과 큰 자리의 변경}
\label{spaces-groupoids-section-bigger-site}

\noindent
처음 읽을 때는 이 절을 건너뛰기를 권한다.
스택의 당김은 Stacks, Section
\ref{stacks-section-inverse-image}에서 정의한다.

\begin{lemma}
\label{spaces-groupoids-lemma-quotient-stack-change-big-site}
큰 자리 $\Sch_{fppf}$와 $\Sch'_{fppf}$가 주어졌다고 하자.
$\Sch_{fppf}$가 $\Sch'_{fppf}$에 포함된다고 가정하자.
Topologies, Section \ref{topologies-section-change-alpha}를 보라.
$S \in \Ob(\Sch_{fppf})$라 하자.
$B, U, R \in \Sh((\Sch/S)_{fppf})$를 대수공간이라 하고,
$(U, R, s, t, c)$를 $B$ 위 대수공간에서의 groupoid라 하자.
$f : (\Sch'/S)_{fppf} \to (\Sch/S)_{fppf}$를 포함 함자
$u : \Sch_{fppf} \to \Sch'_{fppf}$에 대응하는 자리의 사상이라
하자. 그러면 표준적인 동치
$$
[f^{-1}U/f^{-1}R]
\longrightarrow
f^{-1}[U/R]
$$
가 존재하며, 이는 $(\Sch'/S)_{fppf}$ 위 groupoid에서의
스택들의 동치이다.
\end{lemma}

\begin{proof}
Spaces, Lemma \ref{spaces-lemma-change-big-site}에 의해
$f^{-1}B, f^{-1}U, f^{-1}R \in \Sh((\Sch'/S)_{fppf})$는
대수공간이다. 따라서
$(f^{-1}U, f^{-1}R, f^{-1}s, f^{-1}t, f^{-1}c)$는
$f^{-1}B$ 위 대수공간에서의 groupoid이다. 그러므로 명제는
의미가 있다.

\medskip\noindent
범주 $u_p[U/_{\!p}R]$는 범주 $u_{pp}[U/_{\!p}R]$를 사상의
오른쪽 곱셈계 $I$에 관하여 국소화한 것이다.
$u_{pp}[U/_{\!p}R]$의 대상은 삼중항
$$
(T', \phi : T' \to T, x)
$$
이다. 여기서
$T' \in \Ob((\Sch'/S)_{fppf})$,
$T \in \Ob((\Sch/S)_{fppf})$이고, $\phi$는 $S$ 위 스킴의
사상이며 $x : T \to U$는 $(\Sch/S)_{fppf}$ 위 층의 사상이다.
스킴의 사상 $\phi : T' \to T$는 사상
$\phi : T' \to u(T)$와 같은 자료이다. 그리고 $u(T)$가
$f^{-1}T$를 표현하므로 이는 사상 $T' \to f^{-1}T$와 같은
자료이다. 더욱이 대수공간 위의 $f^{-1}$는 완전충실이므로
(Spaces, Lemma \ref{spaces-lemma-fully-faithful} 참조)
$x$를 사상 $x : f^{-1}T \to f^{-1}U$로 생각할 수도 있다.
이제부터는 이러한 동일시를 별도의 언급 없이 사용한다.
사상
$$
(a, a', \alpha) :
(T'_1, \phi_1 : T'_1 \to T_1, x_1)
\longrightarrow
(T'_2, \phi_2 : T'_2 \to T_2, x_2)
$$
은 $u_{pp}[U/_{\!p}R]$에 속하며 가환 도식
$$
\xymatrix{
& & U \\
T'_1 \ar[d]_{a'} \ar[r]_{\phi_1} &
T_1 \ar[d]_a \ar[ru]^{x_1} \ar[r]_\alpha &
R \ar[d]^t \ar[u]_s \\
T'_2 \ar[r]^{\phi_2} &
T_2 \ar[r]^{x_2} &
U
}
$$
이다. 이러한 사상이 $I$의 원소일 필요충분조건은
$T'_1 = T'_2$이고 $a' = \text{id}$인 것이다.
함자
$$
u_{pp}[U/_{\!p}R] \longrightarrow [f^{-1}U/_{\!p}f^{-1}R]
$$
를 대상 위에서 규칙
$$
(T', \phi : T' \to T, x) \longmapsto (x \circ \phi : T' \to f^{-1}U)
$$
으로, 위와 같은 사상 위에서 규칙
$$
(a, a', \alpha) \longmapsto (\alpha \circ \phi_1 : T'_1 \to f^{-1}R)
$$
으로 정의한다.
$I$의 원소들이 동형사상으로 옮겨짐은 명백하다. 실제로
$(f^{-1}U, f^{-1}R, f^{-1}s, f^{-1}t, f^{-1}c)$는
$f^{-1}B$ 위 대수공간에서의 groupoid이다. 따라서 이 함자는
표준적으로 함자
$$
u_p[U/_{\!p}R] \longrightarrow [f^{-1}U/_{\!p}f^{-1}R]
$$
를 통하여 인수분해된다. 스택화를 적용하면 스택의 함자
$$
f^{-1}[U/R] \longrightarrow [f^{-1}U/f^{-1}R]
$$
를 얻으며, 이는 $(\Sch'/S)_{fppf}$ 위의 함자이다.
실제로 Stacks, Lemma
\ref{stacks-lemma-technical-up}에 의해 스택 $f^{-1}[U/R]$는
$u_p[U/_{\!p}R]$의 스택화이다.

\medskip\noindent
이제 스택의 사상을 얻었으며, 이것이 동치임을 확인하려면
완전충실이고 대상들이 국소적으로 본질적 상에 있음을 보이면 충분하다.
Stacks, Lemmas \ref{stacks-lemma-characterize-ff}와
\ref{stacks-lemma-characterize-essentially-surjective-when-ff}를
보라. 대상들에 대한 명제는 $f^{-1}R$가 전사 에탈 사상
$f^{-1}W \to f^{-1}R$를 허용하므로 성립한다. 여기서 $W$는
어떤 $(\Sch/S)_{fppf}$의 대상이다.
% P09-ERR-1466
함자가 ‘충만’임을 보이려면 사상들이 국소적으로 함자의 상에 있음을
보이면 충분하다. 이는 $f^{-1}U$가 전사 에탈 사상
$f^{-1}W \to f^{-1}U$를 허용하므로 성립한다. 여기서 $W$는
어떤 $(\Sch/S)_{fppf}$의 대상이다.
% P09-ERR-1467
함자가 충실하다는 증명은 생략한다.
\end{proof}











\section{분리 조건}
\label{spaces-groupoids-section-separation}

\noindent
이는 실제로 사상 $j : R \to U \times_B U$에 대한 조건을
뜻한다. 이때 대수공간에서의 groupoid $(U, R, s, t, c)$가
$B$ 위에 주어져 있다.
앞 절에서와 같이 먼저 대응하는 도식을 정식화한다.

\begin{lemma}
\label{spaces-groupoids-lemma-diagram-diagonal}
Section \ref{spaces-groupoids-section-notation}에서와 같은 $B \to S$를 잡자.
$(U, R, s, t, c)$를 $B$ 위 대수공간에서의 groupoid라 하자.
$G \to U$를 안정자 군 대수공간이라 하자.
가환 도식
$$
\xymatrix{
R \ar[d]^{\Delta_{R/U \times_B U}} \ar[rrr]_{f \mapsto (f, s(f))} & & &
R \times_{s, U} U \ar[d] \ar[r] & U \ar[d] \\
R \times_{(U \times_B U)} R \ar[rrr]^{(f, g) \mapsto (f, f^{-1} \circ g)} & & &
R \times_{s, U} G \ar[r] & G
}
$$
에서 왼쪽의 두 수평 화살표는 동형사상이고 오른쪽 정사각형은
섬유곱 정사각형이다.
\end{lemma}

\begin{proof}
생략한다. 정의와 대수기하에서의 함자적 관점에 관한 연습문제이다.
\end{proof}

\begin{lemma}
\label{spaces-groupoids-lemma-diagonal}
Section \ref{spaces-groupoids-section-notation}에서와 같은 $B \to S$를 잡자.
$(U, R, s, t, c)$를 $B$ 위 대수공간에서의 groupoid라 하자.
$G \to U$를 안정자 군 대수공간이라 하자.
\begin{enumerate}
\item 다음은 동치이다.
\begin{enumerate}
\item $j : R \to U \times_B U$는 분리 사상이다.
\item $G \to U$는 분리 사상이다.
\item $e : U \to G$는 닫힌 몰입이다.
\end{enumerate}
\item 다음은 동치이다.
\begin{enumerate}
\item $j : R \to U \times_B U$는 국소 분리 사상이다.
\item $G \to U$는 국소 분리 사상이다.
\item $e : U \to G$는 몰입이다.
\end{enumerate}
\item 다음은 동치이다.
\begin{enumerate}
\item $j : R \to U \times_B U$는 준분리 사상이다.
\item $G \to U$는 준분리 사상이다.
\item $e : U \to G$는 준콤팩트 사상이다.
\end{enumerate}
\end{enumerate}
\end{lemma}

\begin{proof}
군 대수공간 $G \to U$는 $R \to U \times_B U$의
기저변환이다. 이 기저변환은 대각사상
$U \to U \times_B U$에 의해 취한다. Lemma
\ref{spaces-groupoids-lemma-groupoid-stabilizer}를 보라. 따라서 $j$가 분리
사상(각각 국소 분리 사상, 준분리 사상)이면 $G \to U$도 분리
사상(각각 국소 분리 사상, 준분리 사상)이다.
Morphisms of Spaces, Lemma
\ref{spaces-morphisms-lemma-base-change-separated}를 보라.
따라서 (1), (2), (3)에서 (a) $\Rightarrow$ (b)이다.

\medskip\noindent
거꾸로 $G \to U$가 분리 사상(각각 국소 분리 사상, 준분리
사상)이면, 사상 $e : U \to G$는 구조 사상 $G \to U$의
단면이므로 닫힌 몰입(각각 몰입, 준콤팩트 사상)이다.
Morphisms of Spaces, Lemma
\ref{spaces-morphisms-lemma-section-immersion}을 보라.
따라서 (1), (2), (3)에서 (b) $\Rightarrow$ (c)이다.

\medskip\noindent
$e$가 닫힌 몰입(각각 몰입, 준콤팩트 사상)이면 Lemma
\ref{spaces-groupoids-lemma-diagram-diagonal}의 결과(그리고 Spaces, Lemma
\ref{spaces-lemma-base-change-immersions}와 Morphisms of Spaces,
Lemma \ref{spaces-morphisms-lemma-base-change-quasi-compact})에 의해
$\Delta_{R/U \times_B U}$가 닫힌 몰입(각각 몰입, 준콤팩트
사상)임을 안다. 따라서 (1), (2), (3)에서
(c) $\Rightarrow$ (a)이다.
\end{proof}
